\documentclass[11pt]{article}
\usepackage[margin=1in]{geometry}
\usepackage{amsmath,amssymb,amsthm,mathtools}
\usepackage{enumitem}
\usepackage{microtype}
\usepackage{hyperref}
\usepackage[T1]{fontenc}
\usepackage{lmodern}
\usepackage{longtable}
\usepackage{array}

\newtheorem{theorem}{Theorem}[section]
\newtheorem{proposition}[theorem]{Proposition}
\newtheorem{lemma}[theorem]{Lemma}
\newtheorem{corollary}[theorem]{Corollary}
\newtheorem{conjecture}[theorem]{Conjecture}
\theoremstyle{definition}
\newtheorem{definition}[theorem]{Definition}
\theoremstyle{remark}
\newtheorem{remark}[theorem]{Remark}

\newcommand{\OP}{\mathcal O}
\newcommand{\Lin}{\operatorname{LinExt}}
\newcommand{\Aut}{\operatorname{Aut}}
\newcommand{\Cyc}{\operatorname{CycExt}}
\newcommand{\Lie}{\operatorname{Lie}}
\newcommand{\sgn}{\operatorname{sgn}}
\newcommand{\Ind}{\operatorname{Ind}}
\newcommand{\ch}{\operatorname{ch}}
\newcommand{\HH}{\mathbb H}
\newcommand{\PLie}{\mathsf{Lie}}
\newcommand{\Sh}{\operatorname{Sh}}
\newcommand{\shuf}{\mathbin{\sqcup\mkern-3mu\sqcup}}
\newcommand{\st}[1]{\textup{[\textsc{#1}]}}
\newcommand{\src}[1]{\nolinkurl{#1}}

\title{The poset Joyal--Klyachko--Stanley theorem:\\
statement, proof, corrections and status}
\author{}
\date{Consolidated write-up --- October 2026}

\begin{document}
\maketitle

\begin{abstract}
This note collects in one place what the repository has established about
the conjecture of \src{notes/poset_jks_research_note.tex} (the \emph{seed
note}). For every finite poset $P$ on $n\ge3$ points, the top homology $H_P$
of the proper part of the lattice of order-preserving partitions is, as an
$\Aut(P)$-module, isomorphic to $\sgn\otimes\PLie(P)$. We give a
self-contained proof via Weisner's theorem, together with a closed formula
for the character. That proof gives an abstract isomorphism. We also record
an explicit isomorphism, given for every $P$ by one choice-free recipe and
commuting with relabelling. We show that naturality along order extensions
holds exactly on connected posets, and that no family of isomorphisms has
it at a disconnected poset on $n\ge3$ points. The seed note's
proposed route through a ``restricted JKS image identity'' is false, and we
record the counterexample. An explicit basis of $H_P$ by splitting cycles of
linear extensions (formerly ``Conjecture~S'') is proved, conditional on a
cited theorem of Wachs. The global module and the count $b_n$ are settled. We
end with a table of computer checks and a list of open problems. The seed
note is not modified. Every statement here carries its epistemic status and
the notes file it comes from.
\end{abstract}

\tableofcontents

\section{How to read this note}

\paragraph{Status tags.} Every result carries one of:
\begin{itemize}[nosep]
\item \st{proven}: a proof is given here or in the cited notes file, and it
  depends only on the inputs it names;
\item \st{computed}: checked by a script in \src{compute/}, listed in
  \S\ref{sec:computations}, with the reproduce command in \src{compute/README.md};
\item \st{cited}: a published theorem, with its location;
\item \st{conjectured}.
\end{itemize}
\st{audited} means that an independent audit re-derived the proof, with its
own code. The audit files are \src{notes/audit_*.md}. A result marked
\st{proven} but not \st{audited} has had only its author's own checks.

\paragraph{Sources.} Each result names the notes file it comes from. The
abbreviations used below are:
\begin{center}
\small
\begin{tabular}{ll}
\textbf{[Seed]} & \src{notes/poset_jks_research_note.tex}\\
\textbf{[Glob]} & \src{notes/global_module_operations_plethlog.md} (Proof I)\\
\textbf{[FPM]} & \src{notes/fixed_point_mobius_identity.md} (Proof II)\\
\textbf{[AudChar]} & \src{notes/audit_character_identity_proof.md}\\
\textbf{[Lit]} & \src{notes/literature_novelty_jks_maps.md}\\
\textbf{[AudLit]} & \src{notes/audit_literature_novelty_jks_maps.md}\\
\textbf{[S]} & \src{notes/conjecture_s_splitting_cycles_span.md}\\
\textbf{[AudS]} & \src{notes/audit_conjecture_s.md}\\
\textbf{[Cone]} & \src{notes/cone_lune_lie_theory.md}\\
\textbf{[AudCone]} & \src{notes/audit_cone_lune_lie_theory.md}\\
\textbf{[Nat]} & \src{notes/natural_isomorphism.md}\\
\textbf{[AudNat]} & \src{notes/audit_natural_isomorphism.md}\\
\textbf{[Bn]} & \src{notes/cyclic_compatible_posets_bn.md}\\
\textbf{[AudBn]} & \src{notes/audit_cyclic_compatible_posets_bn.md}\\
\textbf{[Repro]} & \src{notes/reproduce_character_identity_n_le_5.md}\\
\textbf{[N7]} & \src{notes/n7_character_sweep_and_homology_crosscheck.md}\\
\textbf{[Cit]} & \src{notes/citation_check.md}
\end{tabular}
\end{center}
Equation names such as (mainconj), (characterconj) and (restrictionlemma)
refer to the labels of the seed note.

\paragraph{The seed note's own computations.} The seed note reports
computations for $n\le5$, some $n=6$ examples and values of $b_n$. They are
not used as evidence anywhere below. Wherever a computation is quoted, it is
the repository's own reproduction, with its script.

\section{Objects and conventions}\label{sec:setup}

Let $P$ be a poset on $[n]$. These conventions are those of [Seed], [Repro]
and [AudChar, S1]. Throughout, $n\ge1$, i.e.\ $P$ is nonempty: $e_\zeta$
below involves $\frac1n$, and statements quantified over ``every finite
poset'' are about nonempty posets.

\begin{itemize}
\item \emph{Order-preserving partitions.} A set partition $\pi$ of $[n]$ is
  order-preserving if the quotient digraph $P/\pi$ is acyclic. These form a
  lattice $\OP(P)$ under refinement, with bottom $\hat0$ the discrete
  partition and top $\hat1$ the one-block partition. Write $\widehat{\OP(P)}$
  for its proper part, and
  \[
  H_P=\widetilde H_{n-3}(\Delta(\widehat{\OP(P)});\mathbb C).
  \]
  For a finite bounded poset $L$, $\mu(L)$ means $\mu_L(\hat0,\hat1)$.
\item \emph{Words.} A word $w=w_1\cdots w_n$ is the permutation
  $i\mapsto w_i$. Let $\gamma=(1\,2\,\cdots\,n)$. Then $(w\gamma^j)_i=w_{i+j}$
  (indices mod $n$), so right multiplication rotates positions. For
  $g\in S_n$, $(gw)_i=g(w_i)$, so left multiplication relabels letters.
\item \emph{Cyclic classes.} $\Lin(P)\subseteq S_n$ is the set of linear
  extensions. A cyclic class is a right coset $wC_n$ that meets $\Lin(P)$,
  and $e_C(P)=|\Cyc(P)|$ is the number of them. An element $g\in\Aut(P)$
  fixes a class $C$ with rotation $r=r_g(C)\in\mathbb Z/n$ iff $gw=w\gamma^r$
  for one, equivalently every, $w\in C$. This means $g(w_i)=w_{i+r}$.
\item \emph{The Lie side.} Fix a primitive $n$th root of unity $\zeta$, and
  let $e_\zeta=\frac1n\sum_j\zeta^{-j}\gamma^j$. Then
  \[
  \PLie(P)=\operatorname{span}\{we_\zeta:w\in\Lin(P)\}\subseteq\mathbb C[S_n],
  \]
  with $\Aut(P)$ acting by left multiplication.
\item \emph{Components.} ``Union of components'' means a union of connected
  components of the comparability graph of $P$. For $B\subseteq[n]$, $P|_B$
  is the induced subposet.
\end{itemize}

\paragraph{Cited inputs.} The main theorem uses exactly the following
external results.
\begin{enumerate}[label=(\Alph*)]
\item\label{in:JS} \st{cited} Jen\v{c}a--Sarkoci~\cite{JencaSarkoci},
  Thm~4.5 (arXiv:1112.5782v2 numbering): for $n\ge3$,
  $\Delta(\widehat{\OP(P)})\simeq\bigvee^{e_C(P)}S^{n-3}$. The hypothesis
  is only that $P$ is finite with $n\ge3$; it need not be connected
  ([Lit, \S2.1--2.2]). Their order congruences are the partitions with
  acyclic quotient (their Def.~3.1, after K\"ortesi--Radeleczki--Szil\'agyi;
  equivalently kernels of order-preserving maps, their Thm~3.3, after
  Cz\'edli--Lenkehegyi), and their
  $e_C$ is ours ([Lit, \S2.2]).
\item\label{in:Hall} \st{cited} P.~Hall's theorem: for a finite bounded
  poset $L$, $\mu_L(\hat0,\hat1)=\widetilde\chi(\Delta(\overline L))$, where
  $\overline L$ is the proper part. See Stanley, \emph{EC1}, 2nd ed.,
  Prop.~3.8.5, in the form of Prop.~3.8.6 ([Cit, \S1]).
\item\label{in:Hopf} \st{cited} The Hopf trace formula: for a simplicial
  self-map $f$ of a finite complex,
  $\sum_i(-1)^i\operatorname{tr}(f\mid C_i)=\sum_i(-1)^i\operatorname{tr}(f\mid H_i)$.
  It also holds for the augmented complex and reduced homology. See
  Munkres, \emph{Elements of Algebraic Topology}, Thm~22.1 (number not
  re-checked); an open reference for the augmented form is Wachs,
  \emph{Poset topology: tools and applications}, Thm~2.3.9 ([Cit, \S1]).
\end{enumerate}
Proof~II (Remark~\ref{rem:proofII}) also uses Crapo's complementation
theorem~\cite{Crapo}. Proof~I does not use it.

\section{The theorem}\label{sec:theorem}

\begin{theorem}[Poset JKS theorem]\label{thm:main}
Assume \ref{in:JS}--\ref{in:Hopf}. Let $P$ be a finite poset on $n\ge3$
points. Then for every $g\in\Aut(P)$
\begin{equation}\label{eq:characterconj}
(-1)^{n-3}\mu(\OP(P)^g)=\sgn(g)\sum_{\substack{C\in\Cyc(P)\\gC=C}}\zeta^{r_g(C)},
\end{equation}
which is the seed note's (characterconj). Consequently
\begin{equation}\label{eq:mainiso}
H_P\;\cong\;\sgn\otimes\PLie(P)\qquad\text{as }\Aut(P)\text{-modules}.
\end{equation}
Moreover, if $g$ has order $o$, then
\begin{equation}\label{eq:closed}
\sgn(g)(-1)^{n-3}\mu(\OP(P)^g)\;=\;\sum_{gC=C}\zeta^{r_g(C)}
\;=\;\mu(o)\sum_{x\in F(P,g)}e_C\bigl(P|_{B_0(x)}\bigr).
\end{equation}
Here $\mu(o)$ is the number-theoretic M\"obius function, and $F(P,g)$ is
the set of partitions of $[n]$ into exactly $o$ blocks such that each block
is a union of components and $g$ permutes the blocks cyclically. $B_0(x)$ is
any one block of $x$. $F(P,g)$ is empty unless $\langle g\rangle$ acts
freely on the set of components.
\end{theorem}

\noindent\textbf{Status.} \st{proven}, \st{audited}. Proved independently in
[Glob, \S2, Thm~2.8] (Proof~I) and [FPM] (Proof~II). Both proofs were
audited in [AudChar], with verdict \emph{PROOF HOLDS} for each and no
mathematical error. \st{computed}: \eqref{eq:characterconj} holds for all
unlabeled posets with $n\le7$, and the right-hand equality of
\eqref{eq:closed} has been checked through $n=10$ examples
(\S\ref{sec:computations}).

\begin{remark}[What kind of isomorphism]\label{rem:abstract}
\eqref{eq:mainiso} is an abstract isomorphism: two representations of a
finite group over $\mathbb C$ with the same character are isomorphic. The
seed note's conjecture asks for a \emph{natural} isomorphism. Neither proof
gives one, and neither claims to ([AudChar, \S5]). Choosing an isomorphism
for one poset in each $S_n$-orbit and transporting it by relabelling gives
isomorphisms compatible with relabelling, which is all that
\S\ref{sec:global} needs. Nothing more canonical follows from the character
proof. An explicit, choice-free isomorphism, and the precise extent to which
any isomorphism can be natural, are in \S\ref{sec:natural}.
\end{remark}

The proof below is Proof~I of [Glob, \S2], in the form re-derived in
[AudChar, \S1.1--1.2], with the two expository slips of [AudChar, \S2]
fixed. Slip~(a) is fixed by Lemma~\ref{lem:small}; slip~(b) is fixed in the
proof of Proposition~\ref{prop:Mside}.

\subsection{The two characters}

\begin{lemma}[eq.\ (homchar); {[Glob, Thm~2.8]}, {[FPM, \S1]}, {[AudChar, S5]}]
\label{lem:homchar}
Assume \ref{in:JS}--\ref{in:Hopf} and $n\ge3$. For $g\in\Aut(P)$,
$\operatorname{tr}(g\mid H_P)=(-1)^{n-3}\mu(\OP(P)^g)$. \st{proven}
\end{lemma}

\begin{proof}
$g$ acts on $\OP(P)$ by order automorphisms and so acts simplicially on
$\Delta=\Delta(\widehat{\OP(P)})$. An order automorphism of a finite chain is
the identity. So a chain that $g$ maps to itself is fixed pointwise, with its
orientation. Hence $\operatorname{tr}(g\mid C_i)$ counts the $g$-fixed
$i$-chains, which are the $i$-chains of the proper part of
$L^g=\OP(P)^g$. Note that $\hat0,\hat1\in L^g$. Apply \ref{in:Hopf} on the
augmented complex, with the empty simplex in degree $-1$. Then
$\sum_i(-1)^i\operatorname{tr}(g\mid\widetilde H_i(\Delta))
=\widetilde\chi(\Delta(\overline{L^g}))=\mu(L^g)$ by \ref{in:Hall}. By
\ref{in:JS}, only $i=n-3$ survives.
\end{proof}

\begin{lemma}[eq.\ (liechar); {[AudChar, S1]}]\label{lem:liechar}
Let $P$ be a finite poset on $n\ge1$ points and $g\in\Aut(P)$. Then
$\operatorname{tr}(g\mid\PLie(P))=\sum_{gC=C}\zeta^{r_g(C)}$. \st{proven}
\end{lemma}

\begin{proof}
$\gamma e_\zeta=\frac1n\sum_j\zeta^{-j}\gamma^{j+1}=\zeta e_\zeta$, so
$w\gamma^je_\zeta=\zeta^jwe_\zeta$. The vector $we_\zeta$ is supported on the
coset $wC_n$, with all coefficients nonzero, and distinct cosets are disjoint.
So choosing one $w_C\in C$ for each class gives a basis
$\{w_Ce_\zeta\}_{C\in\Cyc(P)}$ of $\PLie(P)$; in particular
$\dim\PLie(P)=e_C(P)$. Since $g\Lin(P)=\Lin(P)$, $g$ permutes the classes.
A non-fixed class contributes $0$ to the trace. If $gw_C=w_C\gamma^r$, then
$g\cdot w_Ce_\zeta=\zeta^rw_Ce_\zeta$. The rotation $r$ does not depend on
the representative: for $w'=w\gamma^j$, $gw'=w'\gamma^r$.
\end{proof}

\begin{remark}[$\zeta$ conventions]
Both sides of \eqref{eq:closed} are integers. The Lie side reduces to a
Ramanujan sum $\sum_{s\in(\mathbb Z/o)^\times}\omega^s=\mu(o)$, which is the
same for every primitive $\omega$ and under $r\mapsto-r$. So the choice of
$\zeta$ and the direction of rotation do not matter ([AudChar, S2]).
\end{remark}

\subsection{The M\"obius side: Weisner's theorem on the orbit partition}

Fix $g\in\Aut(P)$ of order $o$. Put $L=\OP(P)$, and let $O_g$ be the
partition of $[n]$ into $\langle g\rangle$-orbits. Let $X(P,g)$ be the set of
partitions of $[n]$ whose blocks are unions of components and are permuted
transitively by $g$. Order $X(P,g)$ by refinement; it always contains
$\hat1$.

\begin{lemma}[{[Glob, Lemma~2.1]}]\label{lem:orbitquot}
Let a finite group $K$ act on a poset $R$ by automorphisms. Then the quotient
digraph of $R$ by the $K$-orbits is acyclic. \st{proven}
\end{lemma}

\begin{proof}
Suppose the orbit quotient has a cycle $X_1\to\cdots\to X_r\to X_1$. Lift it
to $a_i<k_ia_{i+1}$ with $a_i\in X_i$, $k_i\in K$ and $a_{r+1}=a_1$. Applying
$k_1$, then $k_1k_2$, and so on, gives
$a_1<k_1a_2<k_1k_2a_3<\cdots<ka_1$ for some $k\in K$. Then
$a_1<ka_1<k^2a_1<\cdots<k^{|K|}a_1=a_1$, a contradiction.
\end{proof}

\begin{lemma}[{[Glob, Lemma~2.2]}, {[AudChar, I2--I4]}]\label{lem:lattice}
\st{proven}
\begin{enumerate}[label=(\roman*)]
\item $L$ is a lattice whose meet is the partition meet.
\item $L^g$ is a lattice containing $\hat0$, $\hat1$ and $O_g$.
\item For $x\in L^g$, the join of $x$ and $O_g$ in $L^g$ is their partition
  join. Its blocks are the $\langle g\rangle$-orbits of blocks of $x$.
\item $\{x\in L^g:x\vee O_g=\hat1\}=X(P,g)$.
\end{enumerate}
\end{lemma}

\begin{proof}
(i) Suppose $P/(x\wedge y)$ had a cycle. Its image in $P/x$ is a closed walk
whose steps are ``$<$'' or ``$=$''. Since $P/x$ is acyclic, every step is
``$=$''. The same holds in $P/y$, so the cycle lies in one block of
$x\wedge y$, a contradiction. A finite meet-semilattice with a top element is
a lattice.

(ii) $g$ acts on $L$ by lattice automorphisms, so $L^g$ is closed under the
meet and the join of $L$. $O_g$ is $g$-fixed, and it is order-preserving by
Lemma~\ref{lem:orbitquot} with $K=\langle g\rangle$ acting on $P$.

(iii) $\langle g\rangle$ acts on the poset $P/x$. The partition join
$x\vee_\Pi O_g$ has as blocks the unions $\bigcup_jg^jB$ of $x$-blocks $B$.
Its quotient is the orbit quotient of $P/x$, which is acyclic by
Lemma~\ref{lem:orbitquot}. So the partition join is in $L^g$. Every upper
bound of $x$ and $O_g$ in $L^g$ lies above it.

(iv) If $x\vee O_g=\hat1$, then $\langle g\rangle$ is transitive on the blocks
of $x$. A poset with a transitive group of automorphisms is an antichain:
$B<g^kB$ would give $B<g^{k|g|}B=B$. So $P/x$ has no relations, and every
block of $x$ is a union of components. Conversely, every $x\in X(P,g)$ is
order-preserving (its quotient is an antichain), is $g$-fixed, and has join
$\hat1$ with $O_g$ by (iii).
\end{proof}

\begin{lemma}[Weisner; {[Glob, Lemma~2.3]}]\label{lem:weisner}
In a finite lattice $M$ with $a\ne\hat0$,
$\sum_{x:\,x\vee a=\hat1}\mu(\hat0,x)=0$. \st{proven}
\end{lemma}

\begin{proof}
Expand $\delta(x\vee a,\hat1)=\sum_{y\ge x\vee a}\mu(y,\hat1)$. Then
\[
\sum_x\mu(\hat0,x)\sum_{y\ge x,\,y\ge a}\mu(y,\hat1)
=\sum_{y\ge a}\mu(y,\hat1)\sum_{x\le y}\mu(\hat0,x)
=\sum_{y\ge a}\mu(y,\hat1)\,\delta(y,\hat0)=0,
\]
since $a\ne\hat0$.
\end{proof}

We apply Weisner's theorem in $M=L^g$ with $a=O_g$. This needs the join of
$L^g$ itself, and Lemma~\ref{lem:lattice}(iii) shows that it equals the
partition join here ([AudChar, I3, \S4]).

\begin{lemma}[{[Glob, Lemma~2.4]}, {[AudChar, I6, S8]}]\label{lem:interval}
Let $x\in X(P,g)$ have $m$ blocks, let $B_0$ be one block, and let
$h=g^m|_{B_0}\in\Aut(P|_{B_0})$. Put $k=|B_0|=n/m$. \st{proven}
\begin{enumerate}[label=(\roman*)]
\item $[\hat0,x]_{L^g}\cong\OP(P|_{B_0})^h$.
\item $(-1)^{n-3}\sgn(g)=(-1)^{k-3}\sgn(h)$.
\end{enumerate}
\end{lemma}

\begin{proof}
(i) There are no relations between blocks of $x$. So $\pi\le x$ is
order-preserving iff each restriction $\pi|_B$ is, and
$[\hat0,x]_L\cong\prod_B\OP(P|_B)$. $g$ permutes the factors cyclically, so a
$g$-fixed element is determined by its $B_0$-component, which must be
$g^m$-fixed. Conversely, an $h$-fixed $\sigma\in\OP(P|_{B_0})$ extends to
$\bigcup_{j<m}g^j\sigma$, which is consistent because $g^m\sigma=\sigma$.

(ii) Each cycle of $g$ meets $B_0$ in exactly one cycle of $h$ and is $m$
times as long. So $g$ and $h$ have the same number $c$ of cycles, and
$\sgn(g)=(-1)^{n-c}$, $\sgn(h)=(-1)^{k-c}$.
\end{proof}

For any finite poset $Q$ on $|Q|\ge1$ points and $h\in\Aut(Q)$, put
\[
\widetilde M(Q,h)=(-1)^{|Q|-3}\sgn(h)\,\mu(\OP(Q)^h).
\]

\begin{lemma}[Base case; fixes slip (a) of {[AudChar, \S2]}]\label{lem:small}
Assume \ref{in:JS}--\ref{in:Hall}. For every finite poset $Q$ on $|Q|\ge1$ points,
$\widetilde M(Q,1)=e_C(Q)$. \st{proven}
\end{lemma}

\begin{proof}
For $|Q|\ge3$, by \ref{in:Hall} and \ref{in:JS},
$\mu(\OP(Q))=\widetilde\chi(\bigvee^{e_C(Q)}S^{|Q|-3})=(-1)^{|Q|-3}e_C(Q)$.
For $|Q|\le2$, \ref{in:JS} says nothing, and we compute by hand
([AudChar, S7]). If $|Q|=1$, then $\OP(Q)$ has one element, so $\mu=1$ and
$\widetilde M=(-1)^{-2}=1=e_C(Q)$. If $|Q|=2$, then whether $Q$ is a chain or
an antichain, $\OP(Q)=\{\hat0<\hat1\}$, so $\mu=-1$ and
$\widetilde M=(-1)^{-1}(-1)=1$. Also $e_C(Q)=1$: the two words on two letters
form one rotation class.
\end{proof}

A homology-free proof of Lemma~\ref{lem:small} exists, in
\src{notes/mu_op_elementary.md}, independently audited in
\src{notes/audit_mu_op_elementary.md}.

\begin{proposition}[M\"obius side; {[Glob, Prop.~2.5]}, {[AudChar, I7--I11]}]
\label{prop:Mside}
Assume \ref{in:JS}--\ref{in:Hall}. For every finite poset $P$ on $n\ge1$ points and every
$g\in\Aut(P)$ of order $o$,
\[
\widetilde M(P,g)=\mu(o)\sum_{x\in F(P,g)}e_C\bigl(P|_{B_0(x)}\bigr).
\]
\st{proven}
\end{proposition}

\begin{proof}
For $x\in X(P,g)$ with $m=m(x)$ blocks, let $B_0(x)$ be one block and
$h_x=g^{m}|_{B_0(x)}$. Different choices of block give isomorphic pairs
$(P|_{B_0},h_x)$, since $g^j$ conjugates one to the other. $\widetilde M$ is
an invariant of such isomorphism classes ([AudChar, I8]).

\emph{The recursion.} Suppose $g\ne1$. Then $O_g\ne\hat0$. By
Lemma~\ref{lem:lattice}(iv), Weisner's theorem in $L^g$ with $a=O_g$ reads
$\sum_{x\in X(P,g)}\mu_{L^g}(\hat0,x)=0$. Multiply by $(-1)^{n-3}\sgn(g)$ and
use Lemma~\ref{lem:interval} term by term:
\begin{equation}\label{eq:W}
\sum_{x\in X(P,g)}\widetilde M\bigl(P|_{B_0(x)},h_x\bigr)=0\qquad(g\ne1).\tag{W}
\end{equation}
This holds for every size $n\ge2$. It uses Lemmas
\ref{lem:orbitquot}--\ref{lem:interval} only, and not \ref{in:JS}.

\emph{M\"obius inversion on $X(P,g)$.} For $x\in X(P,g)$, the map
$y\mapsto\{g^jD:D\in y,\ 0\le j<m(x)\}$ is a bijection from
$X(P|_{B_0(x)},h_x)$ onto $\{z\in X(P,g):z\le x\}$. Its inverse takes the
blocks of $z$ that lie inside $B_0(x)$. $h_x$ is transitive on these because
the stabiliser of $B_0(x)$ in $\langle g\rangle$ is $\langle g^{m(x)}\rangle$.
The bijection matches the pairs (block, induced automorphism) up to
isomorphism: if $y\mapsto z$, then $m(z)=m(x)m(y)$ and
$h_z=h_x^{m(y)}|_{D}$ for the block $D\subseteq B_0(x)$ of $z$
([AudChar, I7]). Put $f(z)=\widetilde M(P|_{B_0(z)},h_z)$. Then for every
$x\in X(P,g)$,
\[
\sum_{z\le x}f(z)=\Phi(x):=[h_x=1]\;e_C\bigl(P|_{B_0(x)}\bigr).
\]
If $h_x\ne1$, this is \eqref{eq:W} for the pair $(P|_{B_0(x)},h_x)$,
transported by the bijection. If $h_x=1$, then $X(P|_{B_0(x)},1)=\{\hat1\}$,
since a trivial group is transitive only on a single block. So
$\{z\le x\}=\{x\}$, and $f(x)=\widetilde M(P|_{B_0(x)},1)=e_C(P|_{B_0(x)})$
by Lemma~\ref{lem:small}. M\"obius inversion on the finite poset $X(P,g)$
gives
\[
\widetilde M(P,g)=f(\hat1)=\sum_{x\in X(P,g)}\mu_X(x,\hat1)\,\Phi(x).
\]

\emph{Which $x$ contribute} (fixes slip (b)). Let $x\in X(P,g)$ have $m$
blocks. The image of $\langle g\rangle$ acting on the blocks of $x$ is
transitive and cyclic, so it is a cyclic group of order $m$ generated by the
image of $g$. Hence $m\mid o$. If $h_x=1$, then $g^m$ fixes $B_0(x)$
pointwise, and so, conjugating by powers of $g$, it fixes every block
pointwise. So $g^m=1$ and $o\mid m$. Hence $h_x=1$ forces $m=o$. Conversely,
$m=o$ gives $g^m=1$ and $h_x=1$. So $\Phi(x)\ne0$ only for
$x\in F(P,g)=\{x\in X(P,g):m(x)=o\}$. If some component $T$ had a nontrivial
stabiliser $g^j$ with $0<j<o$, then every $x\in X(P,g)$ would have
$g^jB=B$ for the block $B\supseteq T$. Then $m(x)\le j<o$, so $F(P,g)$ would
be empty. Hence $F(P,g)=\emptyset$ unless $\langle g\rangle$ acts freely on
the components.

\emph{The M\"obius value.} Let $x\in F(P,g)$. $\langle g\rangle\cong\mathbb Z/o$
acts freely and transitively on the $o$ blocks of $x$. The elements of
$[x,\hat1]_X$ are the coarsenings of $x$ that are $g$-invariant; transitivity
is then automatic. Such a coarsening is determined by the set $H$ of group
elements carrying $B_0(x)$ into its own block, and $H$ is a subgroup: if
$b\in H$, then $H+b$ indexes the block containing $g^bB_0$, which is the
block of $B_0$, so $H+b=H$. So $[x,\hat1]_X$ is isomorphic to the subgroup
lattice of $\mathbb Z/o$, i.e.\ the divisor lattice of $o$, and
$\mu_X(x,\hat1)=\mu(o)$. Summing over $x\in F(P,g)$ proves the claim.
\end{proof}

\subsection{The Lie side: cutting the circle into $g$-translates}

\begin{lemma}[{[FPM, Lemma~6a]}; {[Glob, Lemma~2.6]}]\label{lem:rot}
Let $w\in\Lin(P)$ and $0\le j<n$. Then $w\gamma^j\in\Lin(P)$ iff the prefix
$\{w_1,\dots,w_j\}$ is a union of components. \st{proven}
\end{lemma}

\begin{proof}
Let $A=\{w_1,\dots,w_j\}$. If $A$ is a union of components, no relation joins
$A$ to its complement, so the rotated word is a linear extension. Conversely,
take $a\in A$ and $b\notin A$. In $w$, $a$ precedes $b$, so $b\not<a$. In
$w\gamma^j$, $b$ precedes $a$, so $a\not<b$.
\end{proof}

\begin{proposition}[Lie side; {[Glob, Prop.~2.7]}, {[AudChar, S3--S4]}]
\label{prop:Lside}
For every finite poset $P$ on $n\ge1$ points and every $g\in\Aut(P)$ of order $o$,
\[
\sum_{gC=C}\zeta^{r_g(C)}=\mu(o)\sum_{x\in F(P,g)}e_C\bigl(P|_{B_0(x)}\bigr).
\]
\st{proven}
\end{proposition}

\begin{proof}
If $g=1$, both sides are $e_C(P)$, since $F(P,1)=\{\hat1\}$. Let $g\ne1$, and
fix a reference point $\rho\in[n]$. For $x\in F(P,g)$, let $B_0(x)$ be the
block containing $\rho$.

\emph{Step 0.} Let $C$ be $g$-fixed with rotation $r$. Then $g(w_i)=w_{i+r}$
for $w\in C$, so $g$ is conjugate in $S_n$ to $\gamma^r$, and every cycle of
$g$ has length $n/\gcd(n,r)$. Hence $o=n/\gcd(n,r)$. Put $k=\gcd(n,r)=n/o$.

\emph{Step 1: cut sets are $k$-periodic.} For $w\in C\cap\Lin(P)$, let
$K(w)=\{j\in\mathbb Z/n:w\gamma^j\in\Lin(P)\}$. Then $|K(w)|=|C\cap\Lin(P)|$.
If $w\gamma^j\in\Lin(P)$, then $g(w\gamma^j)=w\gamma^{j+r}\in\Lin(P)$. So
$K(w)$ is invariant under $+r$, hence under the subgroup $k\mathbb Z/n\mathbb Z$
that $r$ generates. In particular $K(w)$ is a union of cosets of
$k\mathbb Z/n\mathbb Z$, each of size $o$.

\emph{Step 2: from $C$ to partitions.} Fix $w\in C\cap\Lin(P)$ and
$s\in K(w)$. For $i\in\mathbb Z/o$ let $A_{s+ik}$ be the set of letters in
positions $s+ik+1,\dots,s+(i+1)k$. The word $w\gamma^{s+ik}$ is a linear
extension whose length-$k$ prefix is $A_{s+ik}$, and $w\gamma^{s+(i+1)k}$ is
a linear extension too. By Lemma~\ref{lem:rot}, $A_{s+ik}$ is a union of
components. Since $g$ shifts positions by $r$, it maps $A_t$ to $A_{t+r}$, so
it permutes the $o$ sets $A_{s+ik}$ cyclically, and transitively because
$r/k$ is a unit mod $o$. So $x_s=\{A_{s+ik}\}_i\in F(P,g)$. Since $o\ge2$, the
arcs are proper, so $x_s=x_{s'}$ iff $s\equiv s'\pmod k$. Call
$(C,x)$ \emph{compatible} if $x=x_s$ for some $s\in K(w)$. Each $g$-fixed
$C$ has exactly $|K(w)|/o\ge1$ compatible partitions.

\emph{Step 3: a bijection.} Given a compatible pair $(C,x)$, read off:
\begin{itemize}[nosep]
\item $u$: the word on the arc that carries $B_0(x)$, read from its starting
  cut. Then $u\in\Lin(P|_{B_0(x)})$.
\item $q\in(\mathbb Z/o)^\times$: the next arc carries $g^qB_0(x)$. Since $g$
  moves arcs forward by $r/k$ steps, $q=(r/k)^{-1}\bmod o$.
\end{itemize}
Conversely, given $x\in F(P,g)$, $u\in\Lin(P|_{B_0(x)})$ and
$q\in(\mathbb Z/o)^\times$, put
$W=u\cdot g^qu\cdot g^{2q}u\cdots g^{(o-1)q}u$, where $g^ju$ is $u$ with $g^j$
applied to each letter. The blocks $g^{jq}B_0$ are disjoint unions of
components that cover $[n]$, and each piece is a linear extension of its
block. So $W\in\Lin(P)$. The map $g$ sends the piece at index $j$ to the
piece at index $j+q^{-1}$, so $gW=W\gamma^{kq^{-1}}$. Its class is
$g$-fixed, with $r=kq^{-1}$, and $(\text{class of }W,x)$ is compatible (take
$s=0$). The two constructions are mutually inverse. Hence compatible pairs
$(C,x)$ correspond bijectively to triples $(x,u,q)$.

\emph{Step 4: the count.} Under the bijection, the cuts of $C$ inside one arc
are the cuts of $u$ in $P|_{B_0}$, by Lemma~\ref{lem:rot}. So
$|K(w)|/o=|K_{B_0}(u)|$, where $K_{B_0}(u)$ is the cut set of $u$ in
$P|_{B_0}$. Let $\omega=\zeta^k$, a primitive $o$th root of unity. Then
$\zeta^r=\omega^{q^{-1}}$, and
\[
\sum_{gC=C}\zeta^{r_g(C)}
=\sum_{(C,x)}\frac{\zeta^{r(C)}}{|K(w)|/o}
=\sum_{x\in F}\;\sum_{u\in\Lin(P|_{B_0(x)})}\frac{1}{|K_{B_0}(u)|}
\sum_{q\in(\mathbb Z/o)^\times}\omega^{q^{-1}}
=\mu(o)\sum_{x\in F}e_C\bigl(P|_{B_0(x)}\bigr).
\]
The middle step uses the fact that each class $C_0$ of $P|_{B_0}$ contains
exactly $|K_{B_0}(u)|$ linear extensions $u$. The last step is the
Ramanujan sum $\sum_{q\in(\mathbb Z/o)^\times}\omega^q=\mu(o)$.
\end{proof}

\subsection{Proof of Theorem~\ref{thm:main}}

By Propositions~\ref{prop:Mside} and~\ref{prop:Lside},
$\sgn(g)(-1)^{n-3}\mu(\OP(P)^g)=\widetilde M(P,g)=\sum_{gC=C}\zeta^{r_g(C)}$,
which is \eqref{eq:closed}. Multiplying by $\sgn(g)$ gives
\eqref{eq:characterconj}. By Lemmas~\ref{lem:homchar} and~\ref{lem:liechar},
the left side of \eqref{eq:characterconj} is the character of $H_P$, and the
right side is the character of $\sgn\otimes\PLie(P)$. Representations of a
finite group over $\mathbb C$ are determined by their characters, which
gives \eqref{eq:mainiso}. \qed

\begin{remark}[Proof II, by Crapo complementation; {[FPM]}, {[AudChar, \S1.3]}]
\label{rem:proofII}
\st{proven} given \ref{in:JS}--\ref{in:Hopf} and Crapo's complementation
theorem~\cite{Crapo} (\st{cited}; theorem number not checked). \st{audited}.
Instead of the full orbit partition, take $z\in L^g$ to merge a single
$\langle g\rangle$-orbit of maximal length. This is order-preserving because
orbits are antichains ([FPM, Lemma~2]). Crapo's theorem expands
$\mu(L^g)$ over the complements $z^\perp$ of $z$ in $L^g$.
\begin{itemize}[nosep]
\item If the cycles of $g$ have unequal lengths, then $z$ has no complement,
  so $\mu(L^g)=0$ ([FPM, Prop.~4]). By Step~0 above, no class is fixed
  either.
\item If $g$ has type $d^k$ with $k\ge2$, then $z^\perp=\{y_T:T\in\mathcal T\}$
  is an antichain. Here $\mathcal T$ is the set of unions of components $T$
  through a fixed point $c_0$ with $P=T\sqcup gT\sqcup\dots\sqcup g^{d-1}T$.
  So $\mu(L^g)=\mu(d)\sum_{T\in\mathcal T}\mu(\OP(T))$ ([FPM, Prop.~5]).
\item The Lie side is a weighted bijection: the number of $g$-fixed classes
  with rotation $r$ is $\sum_{T\in\mathcal T}e_C(T)$ if $\gcd(n,r)=k$, and
  $0$ otherwise ([FPM, Prop.~6]).
\end{itemize}
Proof~I's $F(P,g)$ is $\{y_T\}$, and its $\mu(o)$ is $\mu(d)$, so the two
closed formulas agree ([AudChar, \S3]). The proofs share the Lie-side
convention, eq.~(homchar) and the use of \ref{in:JS}. They are independent
only on the M\"obius side ([AudChar, R4]). In the unequal-length case, the
Bj\"orner--Walker homotopy complementation formula (\st{cited}) also shows
that the proper part of $L^g$ is contractible. The identity does not need
this.
\end{remark}

\subsection{Consequences}

\begin{corollary}[Connected posets; {[Glob, \S3.1]}, {[FPM, Cor.~7]}]
\label{cor:connected}
If $P$ is connected and $n\ge3$, then $\mu(\OP(P)^g)=0$ for all $g\ne1$, and
$H_P\cong\PLie(P)\cong\mathbb C[\Lin P]$ is a free $\Aut(P)$-module of rank
$|\Lin P|/|\Aut P|$. \st{proven}, \st{audited} ([AudChar, II7]).
\end{corollary}

\begin{proof}
A single component cannot be permuted freely by $g\ne1$, so
$F(P,g)=\emptyset$. Every class meets $\Lin(P)$ exactly once by
Lemma~\ref{lem:rot}, and $\Aut(P)$ acts freely on $\Lin(P)$ ($gw=w$ forces
$g=1$). A free module twisted by $\sgn$ is free.
\end{proof}

This covers ordinal sums $P\oplus Q$ of nonempty posets, products of
connected posets, crowns, and substitutions into a connected poset with at
least two elements ([Glob, \S3.1]). Root-of-unity phases occur only for
automorphisms that permute components freely.

\begin{corollary}[Examples; {[FPM, \S6.1--6.2]}]\label{cor:examples}
\st{proven}
\begin{enumerate}[label=(\alph*)]
\item \emph{Antichain.} For $g$ of type $d^k$,
  $\mu(\Pi_n^g)=(-1)^{k-1}\mu(d)\,d^{k-1}(k-1)!$, and $\mu(\Pi_n^g)=0$ for
  every other $g\ne1$. This is the classical fixed-point computation behind
  JKS. Theorem~\ref{thm:main} for the antichain is classical
  JKS~\cite{Stanley82,Hanlon81,Joyal86}, via Klyachko's
  $\Lie_n\cong\Ind_{C_n}^{S_n}\zeta$~\cite{Klyachko}.
\item \emph{Three disjoint $2$-chains, $g$ cycling them.} Here
  $\chi_{H_P}(g)=-1$. This is the seed note's $n=6$ example. It is now
  proved, and it is also \st{computed} by direct homology
  (\S\ref{sec:computations}).
\item \emph{$m$ equal chains of length $\ell$, $g$ an $m$-cycle on them.}
  $\chi_{H_P}(g)=(-1)^{(m-1)\ell}\mu(m)$.
\end{enumerate}
\end{corollary}

\begin{corollary}[$e_C$ over components; {[Glob, \S3.3]}]\label{cor:eC}
Let $K$ be the set of components of $P$. Then
$e_C(P)=\sum_{\pi\in\Pi(K)}(-1)^{|\pi|-1}(|\pi|-1)!\prod_{T\in\pi}|\Lin(P|_T)|$.
\st{proven}; \st{computed} for all $405$ posets with $n\le6$. Not
independently audited.
\end{corollary}

\section{Correction: the restricted image identity is false}
\label{sec:correction}

The seed note, section ``Embedding into the classical JKS module'', fixes a
classical JKS isomorphism
$J:\widetilde H_{n-3}(\overline{\Pi_n})\otimes\sgn\to\Lie_n$ and states that
the conjecture follows from
\begin{equation}\label{eq:RL}
J(H_P\otimes\sgn)=\operatorname{span}\{we_\zeta:w\in\Lin(P)\}=W_Pe_\zeta,
\qquad W_P=\mathbb C\{\Lin(P)\}.\tag{restrictionlemma}
\end{equation}
Its ``Desired restricted JKS lemma'' and its ``second formulation''
$J(H_P\otimes\sgn)=W_P\kappa_n$, with a Klyachko idempotent $\kappa_n$, make
the same claim. Here $H_P\subseteq\widetilde H_{n-3}(\overline{\Pi_n})$ via
the injective inclusion (homembed).

\textbf{All of these are false.} The conjecture itself is unaffected: it is
now Theorem~\ref{thm:main}, proved by characters, not by
\eqref{eq:RL}.

\begin{lemma}[Lemma B; {[Lit, \S3.3]}, {[AudLit, row 1]}]\label{lem:op}
$\OP(P)=\OP(P^{\mathrm{op}})$, so $H_P=H_{P^{\mathrm{op}}}$ as subspaces of
$\widetilde H_{n-3}(\overline{\Pi_n})$. \st{proven}, \st{audited}
\end{lemma}

\begin{proof}
$\pi$ is order-preserving iff $P/\pi$ is acyclic. Reversing every edge
preserves acyclicity.
\end{proof}

\begin{proposition}[Prop.~C; {[Lit, \S3.3]}, {[AudLit, row 2]}]\label{prop:C}
Let $n\ge3$, and let $C$ be the chain $1<2<\cdots<n$. No linear map
$J:\widetilde H_{n-3}(\overline{\Pi_n})\otimes\sgn\to\mathbb C[S_n]$
satisfies \eqref{eq:RL} for both $C$ and $C^{\mathrm{op}}$. \st{proven},
\st{audited}
\end{proposition}

\begin{proof}
$\OP(C)$ is Boolean, so $H_C$ is one-dimensional, and $H_C=H_{C^{\rm op}}$ by
Lemma~\ref{lem:op}. So \eqref{eq:RL} would force
$\mathbb C\,(12\cdots n)e_\zeta=\mathbb C\,(n\cdots21)e_\zeta$. The vector
$we_\zeta$ is nonzero and supported on the coset $wC_n$. Every rotation of
$12\cdots n$ has at most one descent, while $n\cdots21$ has $n-1\ge2$. So
the two cosets are disjoint, and the two lines differ.
\end{proof}

Prop.~C refutes exactly the statement the seed note makes: one $J$, chosen
first, for all $P$ ([AudLit, \S1.1]). Choosing $J$ after $P$ does not help,
as long as $J$ is $S_n$-equivariant. The next proposition uses the splitting
cycles $\rho_w$ and the signed shuffle relations of Prop.~$\Phi$
(\S\ref{sec:basis}).

\begin{proposition}[Prop.~C$'$; {[AudLit, \S1.2]}]\label{prop:Cprime}
Assume the signed shuffle relation $\sum_{u\in1\shuf 2\cdots n}\sgn(u)\rho_u=0$
(part $\Phi(\Leftarrow)$ of Prop.~\ref{prop:Phi}). Let $n\ge3$, $w\in S_n$,
and let $P_w$ be the chain $w_1<\cdots<w_n$. Then no $S_n$-equivariant
linear map $J:\widetilde H_{n-3}(\overline{\Pi_n})\otimes\sgn\to\mathbb C[S_n]e_\zeta$
satisfies $J(H_{P_w}\otimes\sgn)=\mathbb C\,we_\zeta$. \st{proven} given that
relation, which is \st{computed} exactly at chain level for $n\le6$ and
\st{cited} from Wachs for all $n$.
\end{proposition}

\begin{proof}
$H_{P_w}=\mathbb C\rho_w$, so $J(\rho_w\otimes1)=c\,we_\zeta$ with $c\ne0$. Let
$y=\sum_{u\in1\shuf2\cdots n}uw^{-1}$. Then
$y\cdot(\rho_w\otimes1)=\pm\sum_u\sgn(u)\rho_u\otimes1=0$. By equivariance,
$0=c\,y\,we_\zeta=c\,(1\shuf2\cdots n)e_\zeta$. In
$(1\shuf2\cdots n)e_\zeta=\frac1n\sum_u\sum_j\zeta^{-j}u\gamma^j$, the
identity word arises only from $u=12\cdots n$ (with $j=0$) and
$u=23\cdots n1=\gamma$ (with $j=n-1$). Its coefficient is
$(1+\zeta)/n\ne0$, since $\zeta\ne-1$ for $n\ge3$. Contradiction.
\end{proof}

\begin{remark}[Variants; {[AudLit, \S1.2--1.4]}]
\begin{itemize}[nosep]
\item Without the sign twist the same argument works with
  $y=\sum_u\sgn(u)uw^{-1}$; the coefficient becomes
  $(1+(-1)^{n-1}\zeta)/n\ne0$. \st{proven}.
\item For $\kappa_n=\frac1n\sum_\sigma\zeta^{\operatorname{maj}\sigma}\sigma$,
  the fixed-$J$ version fails for all $n\ge3$: $\kappa_n$ and $w_0\kappa_n$
  are never proportional. \st{proven}. The $P$-dependent version needs
  $\Sh_n\kappa_n\ne0$, which is \st{computed} only for $n\le5$.
\item \emph{Not refuted:} an $\Aut(P)$-equivariant $J_P$ with
  $J_P(H_P\otimes\sgn)=W_Pe_\zeta$. Its existence is \emph{equivalent} to
  \eqref{eq:mainiso}, by semisimplicity ([AudLit, \S1.3]), so it is true by
  Theorem~\ref{thm:main}. It does not, however, reduce anything.
\item The obstruction is structural. The annihilator of $\rho_w$ is
  $\Sh_n w^{-1}$, while that of $we_\zeta$ is
  $\mathbb C[S_n](1-e_\zeta)w^{-1}$, and these differ because
  $\Sh_ne_\zeta\ne0$.
\end{itemize}
\end{remark}

\paragraph{What replaces \eqref{eq:RL}.} The correct restricted image
statement uses the target $\mathbb C[S_n]/\Sh_n\cong\Lie_n^*$ in place of
$\mathbb C[S_n]e_\zeta$. It is
$\Phi(H_P\otimes\sgn)=(W_P+\Sh_n)/\Sh_n$ ([Lit, \S3.4, (S1)]), which is
Theorem~\ref{thm:S} below. Composing with right multiplication by the adjoint
Dynkin element $\theta_n^*$, whose kernel is exactly $\Sh_n$ ([Lit, \S3.3]),
gives a version inside $\mathbb C[S_n]$: $J(H_P\otimes\sgn)=W_P\theta_n^*$.

\section{An explicit basis: splitting cycles of linear extensions
(formerly Conjecture S)}\label{sec:basis}

\subsection{Splitting cycles and shuffles}

For $w\in S_n$, the cuts between consecutive letters of $w$ give interval
partitions forming a Boolean lattice $\Pi_w\subseteq\Pi_n$. Its proper part is
an $(n-3)$-sphere. The \emph{splitting cycle} $\rho_w$ is its fundamental
cycle, normalised as in Wachs~\cite{Wachs}: the chain whose $k$-block
partition is $w_1|\cdots|w_{k-1}|\text{rest}$ has coefficient $+1$. Let
$\Sh_n\subseteq\mathbb C[S_n]$ be the span of the proper shuffles $u\shuf v$
($u,v$ nonempty with complementary letter sets), and let
$W_P=\mathbb C\{\Lin(P)\}$.

\begin{proposition}[Prop.~$\Phi$; {[Lit, \S3.2]}, {[AudLit, rows 10--12]}]
\label{prop:Phi}
The map $\rho_\sigma\otimes1\mapsto\sgn(\sigma)\sigma+\Sh_n$ is a
well-defined $S_n$-isomorphism
$\widetilde H_{n-3}(\overline{\Pi_n})\otimes\sgn\to\mathbb C[S_n]/\Sh_n$.
Equivalently, $\sum_\sigma a_\sigma\rho_\sigma=0$ iff
$\sum_\sigma a_\sigma\sgn(\sigma)\sigma\in\Sh_n$.
\end{proposition}

\noindent\textbf{Status.} \st{cited} for general $n$: it is the dual of
Wachs's comb correspondence (Wachs 1998, Lemma~3.1 and Thms~4.2--4.3,
\emph{preprint} numbering) combined with Ree's theorem
$\Lie_n^\perp=\Sh_n$~\cite{Ree,Reutenauer}. The dualisation was re-derived
in [AudLit, \S2.2], and the Wachs statements were checked against the
preprint, not re-proved. \st{computed} and hence \st{proven} outright for
$n\le6$, by exact chain-level computation with no use of Wachs
([AudLit, \S2.1]). As a statement it was not found in this form in the
literature ([Lit, \S3.2]). Write $\Phi(\Leftarrow)$ for well-definedness and
$\Phi(\Rightarrow)$ for injectivity. Given Wachs's basis theorem,
$\Phi(\Rightarrow)$ follows from $\Phi(\Leftarrow)$ by counting dimensions
([AudS, \S C.1]).

\begin{lemma}[Lemma D; {[Lit, \S3.4]}, {[AudLit, row 13]}]\label{lem:D}
For $w\in\Lin(P)$, every partition in $\Pi_w$ is order-preserving for $P$.
Hence $\rho_w\in H_P$. \st{proven}, \st{audited}
\end{lemma}

\begin{proof}
Cutting $w$ into consecutive segments gives the kernel of an isotone
surjection $P\to[k]$, whose quotient is acyclic. Since $\Delta(\widehat{\OP(P)})$
has no $(n-2)$-simplices, $H_P$ is the space of top cycles, and $\rho_w$ is
one.
\end{proof}

\subsection{The theorem}

Fix a component $P_1$ of $P$, the \emph{base component}. For $w\in\Lin(P)$,
call a prefix of $w$ a \emph{union-prefix} if its letter set is a union of
components. The \emph{blocks} of $w$ are the segments between consecutive
union-prefixes. Let $B_P$ be the set of $w\in\Lin(P)$ whose first block
contains $P_1$. Let
\[
R_P=\operatorname{span}\{s\shuf t:\ s\in\Lin(P|_S),\ t\in\Lin(P|_{[n]\setminus S}),\ S\text{ a nonempty proper union of components}\}.
\]
Every word of such an $s\shuf t$ lies in $\Lin(P)$, so $R_P\subseteq W_P\cap\Sh_n$.

\begin{theorem}[{[S, Thms~A, B, C, Cor.~5.5]}]\label{thm:S}
Let $P$ be a finite poset with any base component.
\begin{enumerate}[label=(\alph*)]
\item \emph{(Theorem A.)} $B_P$ is linearly independent modulo $\Sh_n$, and
  $|B_P|=e_C(P)$. For connected $P$, $W_P\cap\Sh_n=0$: no nonzero
  combination of linear extensions of a connected poset is a sum of proper
  shuffles.
\item \emph{(Theorem B.)} $W_P=\operatorname{span}(B_P)\oplus R_P$, and
  $W_P\cap\Sh_n=R_P$, of dimension $|\Lin P|-e_C(P)$.
\item \emph{(Theorem C, formerly Conjecture S.)} If $n\ge3$, then
  $\{\rho_w:w\in B_P\}$ is a basis of $H_P$. In particular
  $\operatorname{span}\{\rho_w:w\in\Lin(P)\}=H_P$. The relations are exactly
  $\sum a_w\rho_w=0\iff\sum a_w\sgn(w)w\in R_P$.
\item \emph{(Cor.~5.5.)} If $n\ge3$, then $H_P\otimes\sgn\cong\mathbb C[\Lin P]/R_P$
  as $\Aut(P)$-modules, via $\rho_w\otimes1\mapsto\sgn(w)w$.
\end{enumerate}
\end{theorem}

\noindent\textbf{Status} (from the independent audit [AudS], verdict
\emph{PROOF HOLDS}, with no gaps):
\begin{itemize}[nosep]
\item (a), (b): \st{proven}, \st{audited}. These are pure word
  combinatorics, with no dependence on Prop.~$\Phi$ or on \ref{in:JS}.
\item (c), basis and spanning: \st{proven} given $\Phi(\Rightarrow)$,
  Lemma~\ref{lem:D} and \ref{in:JS}; \st{audited}. The kernel statement also
  needs $\Phi(\Leftarrow)$ on the generators of $R_P$.
\item (d): \st{proven} given (c) and $\Phi(\Leftarrow)$ on the generators of
  $R_P$; \st{audited}.
\item \st{computed} at chain level, with no use of $\Phi$ or \ref{in:JS},
  for every poset with $3\le n\le6$: (c) and (d) in full, signs and
  equivariance included ([AudS, \S C.2]).
\item If Prop.~$\Phi$ were to fail for some $n\ge7$, then (a) and (b) would
  stand, and (c) and (d) would revert to \st{conjectured} for that $n$
  ([AudS, \S C.1]). $\Phi$ is \st{cited} for $n\ge7$, as stated above.
\end{itemize}

\paragraph{The key tool ([S, Lemma~1]).} For letters $a\ne x$, define
$E_{a,x}:\mathbb C[S_A]\to\mathbb C[S_{A\setminus x}]$ on words by
$E_{a,x}(w)=+(w\setminus x)$ if $x$ immediately precedes $a$,
$-(w\setminus x)$ if $x$ immediately follows $a$, and $0$ otherwise. Then
$E_{a,x}(\Sh_A)\subseteq\Sh_{A\setminus x}$. Indeed, $E_{a,x}$ is the adjoint
of the Lie substitution $a\mapsto[x,a]$; [S] also gives a direct case
analysis that does not use Ree's theorem. \st{proven}, \st{audited}.

\begin{proof}[Proof sketch of (a), following {[S, \S4]}]
Induct on $n$. Let $X=\sum_{w\in B_P}c_ww\in\Sh_n$, and extend $c$ by $0$ to
all words. Take any $w\in\Lin(P)$ with adjacent letters $x=w_i$,
$a=w_{i+1}$, and contract: $Q=P/\{a,x\}$ is a poset, and $E_{a,x}(X)$ lies in
$\Sh$ and is supported on $B_Q$ ([S, Lemmas~2, 5]). By induction
$E_{a,x}(X)=0$. Reading off the coefficient of $w\setminus x$ gives
$c(w)=c(w')$, where $w'$ is $w$ with $x$ and $a$ swapped, if $x\parallel a$;
and $c(w)=0$ if $x<_Pa$. Adjacent incomparable swaps connect $\Lin(P)$, so
$c$ is constant on $\Lin(P)$. It vanishes at some word: a word outside $B_P$
if $P$ is disconnected, or a word with a cover relation adjacent if $P$ is
connected. So $c\equiv0$. Finally, $|B_P|=e_C(P)$ because each cyclic class
consists of the rotations of the block sequence, and exactly one rotation
puts the block containing $P_1$ first ([S, Lemma~4]).
\end{proof}

(b) follows from (a) and a straightening argument modulo $R_P$
([S, Lemma~6]). (c) follows because, by $\Phi(\Rightarrow)$ and (a), the
$e_C(P)$ cycles $\rho_w$, $w\in B_P$, are independent. They lie in $H_P$ by
Lemma~\ref{lem:D}, and $\dim H_P=e_C(P)$ by \ref{in:JS}.

\begin{remark}[Relation to Theorem~\ref{thm:main}]
For connected $P$, Theorem~\ref{thm:S}(c) gives a second, independent proof
of \eqref{eq:mainiso}, with an explicit equivariant isomorphism
$\rho_w\otimes1\mapsto\sgn(w)\,we_\zeta$ ([AudLit, row 19]; [S, \S6]). For
disconnected $P$ it gives the explicit model
$H_P\otimes\sgn\cong\mathbb C[\Lin P]/R_P$. Theorem~\ref{thm:main} then says
that this is abstractly isomorphic to
$\PLie(P)\cong\mathbb C[\Lin P]/(W_P\cap\mathbb C[S_n](1-e_\zeta))$. An
explicit isomorphism between the two, for every $P$, is the truncation map
$\psi_P$ of \S\ref{sec:natural} (Corollary~\ref{cor:psi}). It is
canonical but not natural along extensions, and for disconnected $P$ with
$n\ge3$ no isomorphism is (Theorem~\ref{thm:ext}).
\end{remark}

\begin{remark}[Novelty of Theorem~\ref{thm:S}]
For the antichain, (a) is the classical fact that words beginning with a
fixed letter are a basis of $\mathbb C[S_n]/\Sh_n$ (Radford; Reutenauer,
Ch.~6). For general $P$, no prior statement of (a) or (b) was found
([AudS, \S F]; [S, \S7]). That is not a proof of novelty; the result may be
folklore.
\end{remark}

\begin{remark}[Cone/lune form; {[Cone]}, audited in {[AudCone]}]
[Cone] restates Theorem~\ref{thm:S}(b)--(c) in the language of
Aguiar--Mahajan's top-lunes. Let $\Gamma_{C_P}$ be the span of the chambers
of the order cone $C_P$, and let $\mathrm{Lune}$ be the span of the
non-singleton top-lunes. Then
$\Gamma_{C_P}\cap\mathrm{Lune}=\operatorname{span}\{\text{top-lunes}\subseteq C_P\}$
([Cone, Thm~4.1(a)]; \st{proven}, \st{audited} [AudCone, row~11]). The
lune form of Theorem~\ref{thm:S}(c) ([Cone, Thm~4.1(c)]) is \st{proven}
given Prop.~$\Phi$, Lemma~\ref{lem:D} and \ref{in:JS}, exactly as
Theorem~\ref{thm:S}(c) is ([AudCone, row~12]). [Cone, Lemma~E, Thm~5.2] also
proves, for every finite real central \emph{essential} arrangement, that the
chambers of a \emph{pointed} top-cone are linearly independent modulo
top-lunes (\st{proven}, \st{audited} [AudCone, rows~7--9, \S B.5]). For the braid arrangement this is the
connected case of Theorem~\ref{thm:S}(a), with a geometric proof. The
general-cone form of Conjecture~S is [Cone, Conjecture~G], \st{conjectured}
and \st{computed} for every top-cone of $A_2,A_3,A_4,B_2,B_3$ and on
samples. [Cone] also corrects Aguiar--Mahajan, \emph{Topics}~\cite{AguiarMahajan}, Ex.~10.9:
``non-singleton top-cone'' should read ``non-pointed top-cone''
([AudCone, row~16]). These results are recorded here only as a remark.
\end{remark}

\section{The natural-isomorphism question}\label{sec:natural}

This section is [Nat], as audited in [AudNat] (verdict: Theorem~$\Psi$
\emph{HOLDS}; one defect, in computational evidence only, recorded at the
end of the section). It answers the seed note's request for a
\emph{natural} isomorphism in the following precise form. For every $P$
there is an explicit isomorphism $H_P\otimes\sgn\cong\PLie(P)$ given by one
$P$-uniform recipe with no choices, and it commutes with relabelling
(Theorem~\ref{thm:Theta}). Naturality along order extensions holds exactly on
connected posets, and no family of maps can have it at a disconnected poset
with $n\ge3$ (Theorem~\ref{thm:ext}).

\subsection{Two models and the truncation map}

Let $P$ be a poset on $[n]$. Write $\Sh_n$, $W_P$ and $R_P$ as in
\S\ref{sec:basis}, and put
\[
M_P=W_P/R_P=(W_P+\Sh_n)/\Sh_n,\qquad L_P=\PLie(P)=W_Pe_\zeta .
\]
A \emph{cut} of $w\in\Lin(P)$ is a $k\in\{0,\dots,n-1\}$ such that
$\{w_1,\dots,w_k\}$ is a union of components. Then $w\gamma^k\in\Lin(P)$ iff
$k$ is a cut, so the cyclic class of $w$ is $\{w\gamma^k:k\text{ a cut}\}$
([S, Lemma~4]). Distinct classes lie in distinct rotation orbits, so
$\{w_Ce_\zeta\}$, one $w_C$ per class $C$, is a basis of $L_P$, and
$\dim L_P=e_C(P)=\dim M_P$ by Theorem~\ref{thm:S}(b) ([AudNat, \S1];
\st{proven}, \st{audited}). Let $\mathrm{pr}_P:\mathbb C[S_n]\to W_P$ be the
coordinate truncation to $\Lin(P)$, and define
\[
\psi_P:L_P\to M_P,\qquad y\mapsto\mathrm{pr}_P(y)+\Sh_n .
\]
It is a function of the element $y$, not of a representative, and it lands
in $M_P$ because $W_P\cap\Sh_n=R_P$. It is $\Aut(P)$-equivariant with no sign
twist: $\Aut(P)$ permutes $\Lin(P)$, and $\Sh_n$ is $S_n$-stable.

\subsection{Theorem $\Psi$}

Call $\lambda:[n]\to\mathbb C^\times$ \emph{admissible} for $P$ if
$\lambda([n])=1$ and $\lambda(U)\ne1$ for every nonempty proper union $U$ of
components, where $\lambda(U)=\prod_{x\in U}\lambda(x)$. For $w\in\Lin(P)$ put
\[
s_\lambda(w)=\sum_{k\text{ a cut of }w}\lambda(w_1\cdots w_k)^{-1}\,w\gamma^k .
\]

\begin{lemma}[{[Nat, Lemma~$\Psi$.0]}, {[AudNat, \S2.1]}]\label{lem:psi0}
If $m$ is a cut of $w$, then
$s_\lambda(w\gamma^m)=\lambda(w_1\cdots w_m)\,s_\lambda(w)$. So each cyclic
class $C$ determines a line $\mathbb C\,s_\lambda(C)$. \st{proven},
\st{audited}. The hypothesis $\lambda([n])=1$ is used here.
\end{lemma}

\begin{theorem}[Theorem~$\Psi$; {[Nat, \S3.1]}, {[AudNat, \S2]}]\label{thm:Psi}
Let $\lambda$ be admissible for $P$. Then the class sums $s_\lambda(C)$, one
for each cyclic class $C$ of $P$, are linearly independent modulo $\Sh_n$.
\end{theorem}

\noindent\textbf{Status.} \st{proven}, \st{audited} ([AudNat], rows~1--3,
verdict \emph{HOLDS}, no gap). It is elementary: it uses only the
contraction lemma and Lemmas~2--4 of [S], which are pure word
combinatorics, and Corollary~\ref{cor:psi} adds only Theorem~\ref{thm:S}(b).
Neither uses Prop.~$\Phi$, nor
Lemma~\ref{lem:D}, nor \ref{in:JS} ([AudNat, \S2.4]). \st{computed}
independently for all $1+2+5+16+63+318$ posets with $n\le6$, for every
primitive $\zeta$ and for three random admissible $\lambda$ per poset
([AudNat, \texttt{[psi]}]).

\begin{proof}[Proof sketch, following {[Nat, \S3.1]} and {[AudNat, \S2.2]}]
Induct on $n$. For $n\le2$ it is checked directly; for the $2$-antichain it
uses $\lambda(x)\ne1$. Let $n\ge3$ and $X=\sum_Cc_Cs_\lambda(C)\in\Sh_n$, and
let $\tilde c(v)$ be the coefficient of $v$ in $X$ (zero off $\Lin P$). By
Lemma~\ref{lem:psi0}, $\tilde c(v)=\lambda(v_1\cdots v_m)\tilde c(v\gamma^m)$
at every cut $m$. Contract an adjacent pair $a,x$ to $Q=P/\{a,x\}$ with
weight $\lambda'(a)=\lambda(a)\lambda(x)$ and $\lambda'=\lambda$ elsewhere;
$\lambda'$ is admissible for $Q$. Then $E_{a,x}(X)\in\Sh$, its coefficients
satisfy the twisted relation for $(Q,\lambda')$, and induction gives
$\tilde c(v'[a\to xa])=\tilde c(v'[a\to ax])$ for all $v'\in\Lin(Q)$. As in
Theorem~\ref{thm:S}(a), this makes $\tilde c$ constant on $\Lin(P)$ and zero
if a cover is realised adjacently. If $P$ is disconnected, a proper cut
gives $c=\lambda(U)c$ with $\lambda(U)\ne1$, so $c=0$. If $P$ is connected,
an adjacent cover gives $c=0$.
\end{proof}

Admissibility cannot be dropped. For $\lambda\equiv1$ the class sums are
dependent modulo $\Sh_n$ for \emph{every} disconnected $P$: their sum is
$\sum_{w\in\Lin P}w$, which is a sum of shuffles $s\shuf t$ over a component
and its complement (\st{proven}, \st{audited} [AudNat, \S5]; \st{computed}
for $1,2,6,19,80$ of the $1,2,6,19,80$ disconnected posets, $n=2,\dots,6$).
Roots of unity of non-primitive order also fail on some posets
(\st{computed}, [AudNat, \S2.3]).

\begin{corollary}[{[Nat, Cor.~$\Psi$.1]}, {[AudNat, row~4, \S2.4]}]\label{cor:psi}
For every finite poset $P$ on $n\ge1$ points, $\psi_P:L_P\to M_P$ is an $\Aut(P)$-isomorphism.
\st{proven}, \st{audited}, unconditionally (elementary).
\end{corollary}

\begin{proof}
$\lambda\equiv\zeta$ is admissible, since $\zeta^{|U|}\ne1$ for
$0<|U|<n$. And $\mathrm{pr}_P(we_\zeta)=s_\zeta(w)/n$, because $w\gamma^k\in\Lin P$
iff $k$ is a cut. So $\psi_P$ maps the basis $\{w_Ce_\zeta\}$ to
$\{s_\zeta(C)/n+\Sh_n\}$, which is independent by Theorem~\ref{thm:Psi}.
Both spaces have dimension $e_C(P)$.
\end{proof}

\begin{corollary}[Antichain; {[Nat, Cor.~$\Psi$.2]}, {[AudNat, row~12, \S6]}]\label{cor:klyachko}
$\mathbb C[S_n]e_\zeta\cap\Sh_n=0$ for every primitive $\zeta$. Hence
$y\mapsto y+\Sh_n$ is an $S_n$-isomorphism
$\mathbb C[S_n]e_\zeta\cong\mathbb C[S_n]/\Sh_n$, and $x\mapsto xe_\zeta$ is
injective on $\Lie_n$. With $\Lie_n^\perp=\Sh_n$ (Ree, \st{cited}) this
reproves Klyachko's $\Lie_n\cong\Ind_{C_n}^{S_n}\zeta$~\cite{Klyachko}.
\st{proven} given Ree's theorem, \st{audited}. \emph{Novelty unclear}: no
proof via contractions modulo shuffles was found in one search, but the
antichain case of Theorem~\ref{thm:S}(a) is classical, and this is
plausibly folklore-adjacent ([AudNat, \S6]).
\end{corollary}

\subsection{The explicit isomorphism $\Theta_P$}

\begin{theorem}[{[Nat, Thm~2]}, {[AudNat, rows~5--7]}]\label{thm:Theta}
For every poset $P$ with $n\ge3$,
\[
\Theta_P=\psi_P^{-1}\circ\Phi_P:\ H_P\otimes\sgn\to\PLie(P),
\qquad \Phi_P(\rho_w\otimes1)=\sgn(w)\,w+R_P,
\]
is an $\Aut(P)$-isomorphism, and $\Theta_{\sigma P}(\sigma h)=\sigma\,\Theta_P(h)$
for every relabelling $\sigma\in S_n$. Its inverse is
\[
\Theta_P^{-1}(we_\zeta)=\frac{\sgn(w)}{n}\sum_{k\text{ a cut of }w}(-1)^{(n-1)k}\zeta^{-k}\,\rho_{w\gamma^k}\otimes1 .
\]
In words: a cyclic class goes to the $\zeta$-twisted signed sum of the
splitting cycles of its members. On connected $P$ this is
$we_\zeta\mapsto\frac{\sgn(w)}n\rho_w\otimes1$.
\end{theorem}

\noindent\textbf{Status.} \st{proven} given the inputs of
Theorem~\ref{thm:S}(d) (Prop.~$\Phi$, Lemma~\ref{lem:D}, \ref{in:JS}),
\st{audited}. Those inputs enter only through
$\Phi_P:H_P\otimes\sgn\cong M_P$, which is Theorem~\ref{thm:S}(d); the
isomorphism $L_P\cong M_P$ (Corollary~\ref{cor:psi}) needs none of them. The
sign comes from $\sgn(\gamma)=(-1)^{n-1}$; the twist lives entirely in
$\Phi_P$, since $\psi_P$ has none ([AudNat, \S2.4]).

\begin{corollary}[Third proof of \eqref{eq:mainiso}]\label{cor:thirdproof}
Theorem~\ref{thm:Theta} gives $H_P\cong\sgn\otimes\PLie(P)$ as
$\Aut(P)$-modules for every $P$ with $n\ge3$, with an explicit isomorphism and
without the character identity. Its inputs are those of
Theorem~\ref{thm:S}(c)--(d), not those of \S\ref{sec:theorem}. \st{proven}
given Prop.~$\Phi$, Lemma~\ref{lem:D} and \ref{in:JS}; \st{audited}
([AudNat, row~6]).
\end{corollary}

\begin{remark}[In what sense $\Theta$ is natural]\label{rem:Thetasense}
This is the precise sense audited in [AudNat, \S3.1].
\begin{itemize}[nosep]
\item $\Theta$ is a natural isomorphism on the groupoid of posets on $[n]$
  and relabellings. On its own this is contentless: any choice of one
  $\Aut$-isomorphism per isomorphism class, transported, is natural on that
  groupoid (Remark~\ref{rem:abstract}).
\item The content is that $\Theta_P$ is given by one recipe, uniform in $P$,
  with no choice of representative, base component or basis: truncate to
  $\Lin(P)$, read modulo $\Sh_n$, apply $\Phi_P^{-1}$. ``Canonical'' is meant
  in this informal sense, not a functorial one. The recipe depends on
  $\zeta$, exactly as $\PLie(P)$ does.
\item $\Theta$ is \emph{not} natural along order extensions once
  disconnected posets are present (Theorem~\ref{thm:ext}(i)). Concretely, it
  fails along some arrow $P\to C_w$ to a linear extension for every
  disconnected $P$ with $3\le n\le6$ (\st{computed}, $2/2$, $6/6$, $19/19$,
  $80/80$, [AudNat, \texttt{[prc]}]), and for every disconnected $P$ with
  $n\ge3$ by Theorem~\ref{thm:ext}(i).
\item There is no conflict with Props.~\ref{prop:C} and~\ref{prop:Cprime}.
  The family $\{\Theta_P\}$ is not the restriction of one ambient map: its
  $P$-dependence is exactly the truncation $\mathrm{pr}_P$. Each $\Theta_P$
  is only $\Aut(P)$-equivariant, which is the case that Prop.~\ref{prop:Cprime}
  does not exclude ([AudNat, \S3.2]). The obvious uniform candidate
  $y\mapsto y+\Sh_n$ carries $L_P$ onto $M_P$ for only $4/5$, $6/16$, $8/63$,
  $10/318$ posets ($n=3,\dots,6$), and never for a chain (\st{computed},
  [AudNat, \texttt{[prc]}]).
\end{itemize}
\end{remark}

\subsection{Naturality along extensions holds exactly on connected posets}

Let $\mathcal E$ be the category of posets on $[n]$ with an arrow $P\to P'$
whenever $\sigma(P)\subseteq P'$ for some $\sigma\in S_n$. Both
$P\mapsto H_P\otimes\sgn$ and $P\mapsto\PLie(P)$ are contravariant on
$\mathcal E$, by $H_{P'}\subseteq H_{\sigma P}$ and
$W_{P'}e_\zeta\subseteq W_{\sigma P}e_\zeta$. For $w\in\Lin(P)$ let $C_w$ be
the chain $w_1<\dots<w_n$.

\begin{theorem}[{[Nat, Thm~1]}, {[AudNat, rows~9--11, \S4]}]\label{thm:ext}
Let $P$ be a poset on $[n]$.
\begin{enumerate}[label=(\roman*)]
\item If $P$ is disconnected and $n\ge3$, then every linear map
  $\eta_P:H_P\otimes\sgn\to\PLie(P)$ with
  $\eta_P(H_{C_w}\otimes\sgn)\subseteq\mathbb Cwe_\zeta$ for all
  $w\in\Lin(P)$ is \emph{not injective}. So no family of isomorphisms is
  natural along the arrows $P\to C_w$ of $\mathcal E$.
\item If $P$ is connected, $\eta_P(\rho_w\otimes1)=\sgn(w)\,we_\zeta$ is an
  $\Aut(P)$-isomorphism. On the full subcategory of connected posets these
  maps form a natural isomorphism, and every natural isomorphism there is
  $c_n\eta$ for one scalar $c_n$ per $n$.
\item $R_Pe_\zeta=0$ if and only if $P$ is connected or $n\le2$.
\end{enumerate}
\end{theorem}

\noindent\textbf{Status.} (i): \st{proven} given Prop.~$\Phi$ and
Lemma~\ref{lem:D}, \st{audited}. The proof takes a nonempty proper union $S$
of components with $|[n]\setminus S|\ge2$, $s\in\Lin(P|_S)$,
$t\in\Lin(P|_{[n]\setminus S})$, and the word $u=t_1st_2\cdots t_b$; no
proper rotation of $u$ lies in $s\shuf t$, so applying $\eta_P$ to the
relation $\sum_{z\in s\shuf t}\sgn(z)\rho_z=0$ forces $\eta_P(\rho_u\otimes1)=0$,
while $\rho_u\ne0$. (ii): \st{proven}, \st{audited}; for $n\ge3$, that $\eta_P$ is an
isomorphism is Theorem~\ref{thm:S}(d) together with $R_P=0$ for connected
$P$, so it rests on the same inputs. (iii): \st{proven}, \st{audited},
elementary. The clause ``or $n\le2$'' is needed: for the $2$-antichain,
$(12+21)e_{-1}=0$. \st{computed} for all posets with $n\le6$ at two primes,
the witness coefficient being exactly $1/n$ in all $8458$ cases
$(P,S,s,t)$ with $3\le n\le6$ ([AudNat, \S4.2]).

Prop.~\ref{prop:C} is the case $P=A_n$ of (i). Theorem~\ref{thm:ext}(i) uses
only the arrows from $P$ to its own linear extensions, and needs no
relabelling naturality.

\subsection{Sum structures}

\paragraph{Ordinal sums ([Nat, Prop.~3]; [AudNat, row~13]).} Write $P\oplus Q$
for all of $P$ below all of $Q$, so $\Lin(P\oplus Q)=\Lin(P)\Lin(Q)$. The
structure maps are deconcatenation $q^M:M_{P\oplus Q}\to M_P\otimes M_Q$,
$uv\mapsto\bar u\otimes\bar v$, and
$q^L:\PLie(P\oplus Q)\to\PLie(P)\otimes\PLie(Q)$,
$uve_\zeta\mapsto ue_{\zeta_P}\otimes ve_{\zeta_Q}$, with $\zeta_P,\zeta_Q$
primitive roots of orders $|P|,|Q|$ (the audit's notational correction).
\begin{itemize}[nosep]
\item Families of $\Aut$-isomorphisms compatible with $q^M,q^L$ exist, but
  they are built from arbitrary choices on the ordinally indecomposable
  pieces, so compatibility with ordinal sums singles out no isomorphism.
\item No compatible family is $c\cdot(w\mapsto we_\zeta)$ on all connected
  posets. In particular $\Theta$, and any family natural along extensions
  among connected posets, is not compatible with ordinal sums. (Take $R=P\oplus1$
  with $P$ disconnected, $|P|\ge3$; compatibility would need
  $R_Pe_\zeta=0$, against Theorem~\ref{thm:ext}(iii).)
\end{itemize}
\st{proven}, \st{audited} by re-reading only (not recomputed).

\paragraph{Disjoint sums ([Nat, Prop.~4]; [AudNat, row~14, \S4.3]).} Let $P$
be on $S$ and $Q$ on $T=[n]\setminus S$. On the $e_\zeta$ side, the
restriction $w\mapsto w|_S\otimes w|_T$ and the deconcatenation never descend
to maps $\PLie(P+Q)\to\PLie(P)\otimes\PLie(Q)$, and the cobracket descends iff
$P$ and $Q$ are both connected and $|S|=|T|$. On the $M$ side the cobracket
$M_{P+Q}\to M_P\otimes M_Q$ always descends. \st{proven}, \st{audited} by
re-reading, including the tensor-intersection step that [Nat] leaves
implicit; \st{computed} in [Nat] for all $71$ pairs with $n\le5$ (not re-run
by [AudNat]). So disjoint-sum compatibility is posed on the $M$ side: on the
$e_\zeta$ side it exists only after transport by an isomorphism such as
$\Theta$, where it holds by definition.

\paragraph{The component-sum map $\varphi_\Sigma$ ([Nat, \S3.3]; [AudNat, rows~15--16]).}
For a component $P_i$, let $B_i$ be the words of $\Lin(P)$ whose first block
contains $P_i$ (a basis of $M_P$, Theorem~\ref{thm:S}(a)--(b)), and let
$\varphi_i:M_P\to\PLie(P)$, $b\mapsto be_\zeta$ for $b\in B_i$. Then
$\varphi_\Sigma=\sum_i\varphi_i$ is $\Aut(P)$-equivariant and involves no
choice (\st{proven} in [Nat]; equivariance not re-checked by [AudNat]). It is invertible for every disconnected $P$ with
$n\le6$ (\st{computed} in [Nat]; re-\st{computed} in [AudNat] for $n\le5$
only). For the antichain, $\varphi_\Sigma(\bar x)=x\theta_n^*e_\zeta$ with
$\theta_n$ the Dynkin element (\st{proven} in [Nat]; not re-checked by
[AudNat]).

\begin{conjecture}[{[Nat, Conj.~3.3]}]\label{conj:phisigma}
$\varphi_\Sigma$ is invertible for every finite poset $P$ on $n\ge1$ points. \st{conjectured}.
\end{conjecture}

\paragraph{The functorial model.} The model $M_P\subseteq\mathbb C[S_n]/\Sh_n$
is natural in every sense above: it is the restriction of the single
$S_n$-isomorphism of Prop.~\ref{prop:Phi}, and it carries injections along
extensions, deconcatenation along ordinal sums and the cobracket along
disjoint sums ([Nat, Prop.~5]; \st{proven} given Prop.~$\Phi$; read in
[AudNat, \S4.3]). So the natural object is $W_P$ modulo shuffles. The
$e_\zeta$ model is natural along extensions exactly on connected posets, and
for every $P$ it is identified with the natural object canonically, but not
functorially, by $\psi_P$.

\paragraph{Defect found by the audit.} The \texttt{[wgt]} check of
\src{compute/natural_isomorphism.py} takes as class representative the
least rotation of a word, which need not lie in $\Lin(P)$, so on such posets
it does not compute $s_\lambda(C)$ ([AudNat, \S5]). Its counts are not
evidence for Theorem~\ref{thm:Psi}, and its $\lambda\equiv1$ counts
($5/6$, $16/19$, $63/80$ at $n=4,5,6$) are wrong: the true counts are
$6/6$, $19/19$, $80/80$, as proven above. The statement for random admissible
weights rests instead on [AudNat, \texttt{[psi]}]. Theorem~\ref{thm:Psi}
itself is unaffected, and so is the \texttt{[psi]} check of [Nat].
\emph{Fixed} in commit \texttt{7c10288} (mg-dab53): the check now takes a
representative in $\Lin(P)$, prints $6/6$, $19/19$, $80/80$, and fails unless
every disconnected $P$ drops.

\section{The global module and $b_n$}\label{sec:global}

Let $\HH_n=\bigoplus_{P\text{ on }[n]}H_P$ and
$V_n=\bigoplus_{P\text{ on }[n]}\PLie(P)$, both made into $S_n$-modules by
relabelling. Let $B_n$ be the set of labelled posets on $[n]$ having some
rotation $L_c=(c+1,\dots,n,1,\dots,c)$ of $12\cdots n$ as a linear
extension, and $b_n=|B_n|$ ([Bn, \S1]). $C_n=\langle\gamma\rangle$ acts on
$B_n$ by relabelling. Let $a_n$ be the number of naturally labelled posets on
$[n]$, and $A(x)=1+\sum_{n\ge1}a_nx^n$.

\begin{proposition}[{[AudBn, \S A, \S C]}]\label{prop:Ind}
$V_n\cong\Ind_{C_n}^{S_n}(\zeta\otimes\mathbb C[B_n])$ as $S_n$-modules.
\st{proven}, unconditionally; \st{audited} by a second read ([AudBn, \S C]).
\end{proposition}

\begin{proof}
By Lemma~\ref{lem:liechar}, choosing a representative $w_O$ in each right
coset $O=wC_n$ makes $\{w_Oe_\zeta\}$ a basis of $\PLie(P)$, indexed by the
$O$ with $O\cap\Lin(P)\ne\emptyset$. For $\tau\in S_n$ we have
$\tau w_O=w_{\tau O}\gamma^j$ for some $j$, so $\tau$ maps the basis vector of
$(P,O)$ to $\zeta^j$ times that of $(\tau P,\tau O)$. So $V_n$ is monomial on
$X=\{(P,O):O\cap\Lin(P)\ne\emptyset\}$. The projection $(P,O)\mapsto O$ onto
$S_n/C_n$ is equivariant and the base is transitive, so $V_n$ is induced from
the stabiliser $C_n$ of $O_0=C_n$ acting on its fibre. The fibre is
$\{P:\text{some }\gamma^j=L_j\in\Lin(P)\}=B_n$. With $w_{O_0}=\mathrm{id}$,
$\gamma^j$ acts on it by $\zeta^j$ times relabelling.
\end{proof}

\begin{corollary}[eq.\ (globalind); {[Glob, \S1(a)]}, {[AudChar, R3]}]
\label{cor:globalind}
Assume \ref{in:JS}--\ref{in:Hopf}. Then
$\HH_n\cong\sgn\otimes V_n\cong\sgn\otimes\Ind_{C_n}^{S_n}(\zeta\otimes\mathbb C[B_n])$,
which is the seed note's (globalind). \st{proven}, \st{audited}.
\end{corollary}

\begin{proof}
On both sides the summands over one $S_n$-orbit of posets form
$\Ind_{\Aut P}^{S_n}$ of a single summand. Apply \eqref{eq:mainiso} for one
$P$ in each orbit, and use
$\Ind_G(\sgn|_G\otimes M)\cong\sgn\otimes\Ind_GM$. Only $\Aut(P)$-equivariance
is used, not naturality ([AudBn, \S C]).
\end{proof}

\begin{theorem}[{[Bn, \S3]}, {[AudBn]}]\label{thm:bn}
$b_n=n[x^n]\log A(x)$; equivalently $na_n=\sum_{k=1}^nb_ka_{n-k}$.
\st{proven}, \st{audited}.
\end{theorem}

\begin{proof}[Proof sketch]
For $P\in B_n$, let $C(P)=\{c:L_c\in\Lin(P)\}\ne\emptyset$. The arcs of the
circle $1\to2\to\cdots\to n\to1$ between consecutive points of $C(P)$ are
mutually incomparable, and $P$ restricted to each arc is irreducible with
respect to the arc order (it does not split as an incomparable prefix and
suffix). This gives a bijection between $B_n$ and the \emph{cyclic
factorizations} of size $n$: a nonempty set of cuts on the circle together
with an irreducible poset on each arc ([Bn, Thm]). By the unique
factorization of naturally labelled posets into irreducibles,
$A=1/(1-I)$, where $I$ counts irreducibles ([Bn, Lemma~2]). Hence cyclic
factorizations are counted by $n[x^n]\log\frac1{1-I}=n[x^n]\log A$
([Bn, Lemma~3]).
\end{proof}

\begin{proposition}[{[Bn, Props.~1--3]}, {[Glob, \S1(b)]}]\label{prop:frob}
\st{proven}. Parts (a)--(c) are \st{audited} in [AudBn]. Part (d) applies
$\omega$ to (c) ([Glob, \S1(b)]); it is elementary and \st{computed}, but it
has had no separate audit.
\begin{enumerate}[label=(\alph*)]
\item $\sum_{P\text{ on }[n]}e_C(P)=(n-1)!\,b_n=\dim V_n$.
\item $\#\mathrm{Fix}(\gamma^k\text{ on }B_n)=b_{\gcd(k,n)}$. The antichain
  is the only $\gamma$-fixed poset, and $b_p\equiv1\pmod p$ for $p$ prime.
\item $\ch V_n=\frac1n\sum_{d\mid n}\mu(d)\,b_{n/d}\,p_d^{n/d}$ and
  $\sum_n\ch V_n=\sum_{d\ge1}\frac{\mu(d)}d\log A(p_d)$. This is the seed
  note's (plethlog). It is the \emph{untwisted} series.
\item With Corollary~\ref{cor:globalind}, the outer twist is
  $\otimes\,\sgn$, i.e.\ $\omega$ on symmetric functions, and
  \[
  \ch\HH_n=\frac1n\sum_{d\mid n}(-1)^{n-n/d}\mu(d)\,b_{n/d}\,p_d^{n/d},
  \]
  which is the seed note's (frob) exactly as printed. The twisted generating
  function is
  $\sum_n\ch\HH_n=\sum_d\frac{\mu(d)}d\log A\bigl((-1)^{d-1}p_d\bigr)$.
\end{enumerate}
\end{proposition}

The sign in (d) is \st{computed} by direct homology, with no use of
Theorem~\ref{thm:main}: the character of $\HH_n$ equals the twisted (frob) on
every class of $S_n$ for $n\le6$. At $n=6$ the twist is decisive. On class
$(6)$, $\HH_6$ has $-1$ where the untwisted formula gives $+1$; on class
$(2,2,2)$ it has $+128$ where the untwisted formula gives $-128$
([Glob, \S1(c)]).

\begin{corollary}[Prime degree]\label{cor:prime}
For $p$ an odd prime,
$\HH_p\cong\frac{b_p-1}{p}\,\mathbb C[S_p]\oplus(\sgn\otimes\Lie_p)$. For
example, $\HH_3\cong5\operatorname{Reg}\oplus(\sgn\otimes\Lie_3)$ and
$\HH_5\cong312\operatorname{Reg}\oplus(\sgn\otimes\Lie_5)$. \st{proven},
given \ref{in:JS}--\ref{in:Hopf}. The same formula holds for $p=2$, which
lies outside the range $n\ge3$ of Theorem~\ref{thm:main}, by direct check:
$\HH_2\cong\mathrm{triv}\oplus\mathbb C[S_2]$, with the convention
$H_P=\widetilde H_{-1}(\emptyset)=\mathbb C$ (\st{proven} by hand in \src{notes/audit_poset_jks_theorem.md}, \S2).
\end{corollary}

\begin{proof}
By Prop.~\ref{prop:frob}(b), $B_p$ consists of the antichain, which is
fixed, and $(b_p-1)/p$ free $C_p$-orbits. A free orbit contributes
$\Ind_{C_p}^{S_p}(\zeta\otimes\mathbb C[C_p])\cong\mathbb C[S_p]$, and the
antichain contributes $\Ind_{C_p}^{S_p}\zeta\cong\Lie_p$~\cite{Klyachko}.
Twist by $\sgn$ (Cor.~\ref{cor:globalind}), and note that
$\sgn\otimes\mathbb C[S_p]\cong\mathbb C[S_p]$. The values $b_3=16$ and
$b_5=1561$ are \st{computed} ([Bn]).
\end{proof}

[Bn, \S4] had recorded this decomposition as conditional on (globalind);
Corollary~\ref{cor:globalind} removes the condition. The deduction is made
here, from the cited results.

\begin{remark}[Why a plethystic logarithm; {[Glob, \S4]}; \textbf{not
independently audited}]
Let $\mathbf Q$ be the species of pairs $(P,\ell)$ with $\ell\in\Lin(P)$, with
product given by disjoint union and concatenation, and coproduct given by
restriction. It is a connected cocommutative Hopf monoid, and $S_n$ acts
freely on $\mathbf Q[n]$, so its cycle index is $A(p_1)$.
Cartier--Milnor--Moore and PBW for species (\st{cited}: Aguiar--Mahajan,
\emph{Monoidal Functors, Species and Hopf Algebras}, Ch.~11; theorem
numbers not re-checked) give
$\ch\operatorname{Prim}(\mathbf Q)=\operatorname{Log}A(p_1)$. So
$\operatorname{Prim}(\mathbf Q)[n]\cong V_n$ as $S_n$-modules. This explains
the shape of (plethlog), and it gives $V_n$ a Lie bracket, the commutator
in $\mathbf Q$. The isomorphism is not graded by posets: the primitives are
not $P$-homogeneous already at $n=2$. \st{proven} modulo the cited theorems,
and \st{computed} for $n\le4$, in [Glob]. The audit [AudChar] covered only
[Glob, \S2]; this part has had no independent audit.
\end{remark}

\section{Computational record}\label{sec:computations}

All scripts are plain Python~3 in \src{compute/}. They are run from that
directory, and archived outputs are in \src{compute/output/}. Each script
exits with status $0$ iff its checks pass, and each has positive controls
(deliberately wrong variants that must fail). The exact commands are in
\src{compute/README.md}. Ranks are taken modulo a large prime, and they are
used only in directions that are valid over $\mathbb Q$ (rank$_p\le$
rank$_{\mathbb Q}$, or a sandwich between two equal bounds).

{\small
\begin{longtable}{>{\raggedright}p{0.36\textwidth}>{\raggedright}p{0.22\textwidth}>{\raggedright}p{0.17\textwidth}p{0.13\textwidth}}
\textbf{What was checked} & \textbf{Script} & \textbf{Range} & \textbf{Source}\\\hline
\endhead
Character identity \eqref{eq:characterconj}, exact in $\mathbb Z[\zeta]$; trace on $\PLie(P)$ computed in $\mathbb C[S_n]$; seed note's $n=6$ examples &
\src{poset_jks_check.py} & all posets $n\le6$ ($758$ pairs at $n=6$) & [Repro]\\
Character identity at $n=7$ & \src{n7_sweep.py} & all $2045$ posets, $4804$ pairs & [N7]\\
Trace of $g$ on $\ker\partial_{n-3}$ by direct homology vs.\ $(-1)^{n-3}\mu(\OP(P)^g)$ (no M\"obius, no Lefschetz) &
\src{equivariant_homology.py} & every pair $n=4,5,6,7$ ($n=4,5$ not archived) & [N7]\\
Steps of Proof I (T1--T5); $\HH_n$ vs.\ (frob) by direct homology; primitives of $\mathbf Q$; $e_\zeta$-bracket test &
\src{poset_lie_structure.py} & $n\le6$ (+ $n=7,8$ examples); $\mathbf Q$: $n\le4$; bracket: $n\le5$ & [Glob]\\
Every lemma of Proof II &
\src{fixed_point_mobius.py} & all posets $n\le6$; $194$ families at $n=7,8,9$ & [FPM]\\
Independent audit of both proofs (checks A--F, H: closed formula, Weisner, Crapo, direct homology) &
\src{audit_character_identity.py} & all posets $n\le6$; $19$ families at $n=7,8,9$; one family at $n=10$ & [AudChar]\\
Prop.~$\Phi$ (chain level), Prop.~C, $\kappa_n$, Conjecture~S ranks, Lemma~D &
\src{splitting_cycles.py} & $\Phi$: $n\le5$; S: $n\le6$ & [Lit]\\
Audit: $\Phi$ at chain level; Conjecture~S at chain level; $\mu(\OP(P)^g)=0$ for connected $P$; Prop.~C$'$ coefficient; Ex.~10.9 &
\src{audit_jks_maps.py} & $n\le6$ (Prop.~C, $\kappa_n$, Ex.~10.9: $n\le5$) & [AudLit]\\
Contraction lemma; $W_P\cap\Sh_n=R_P$; first-block basis; $\rho$-independence &
\src{conjecture_s.py} & $n\le6$ archived; $n\le7$ run, not archived; $\rho$: $n\le5$ & [S]\\
Audit of Theorem~\ref{thm:S}: all lemmas; Thms A, B for every base component; chain-level (c), (d) without $\Phi$ or \ref{in:JS} &
\src{audit_conjecture_s.py} & $n\le6$ & [AudS]\\
$a_n$, $b_n=|B_n|$, $b_n=n[x^n]\log A$, $\sum e_C=(n-1)!b_n$, $C_n$-action on $B_n$ &
\src{cyclic_posets.py} & $n\le7$ ($a_8$, $b_8$ via Lemma~1(c)) & [Bn]\\
Independent audit of the above; character of $V_n$ &
\src{audit_cyclic_posets.py} & $n\le6$ (\texttt{--b8}: $n\le8$) & [AudBn]\\
$V_n\cong\Ind_{C_n}^{S_n}(\zeta\otimes\mathbb C[B_n])$ by direct linear algebra in $\mathbb C[S_n]$ &
\src{second_read_globalind.py} & $n\le5$ & [AudBn, \S C]\\
Lune form of Conjecture~S; pointed cones; Conjecture~G on other arrangements &
\src{cone_lie.py} & $A_2,A_3,A_4,B_2,B_3$; samples of $D_4,B_4$ & [Cone]\\
Independent audit of [Cone]: Zaslavsky face counts, Lemmas L and E exactly, Thm~5.2 (pointed cones), Ex.~10.9 correction, Conjecture~G1 &
\src{audit_cone_lune.py} & $6$ arrangements outside [Cone]'s list ($A_1\times A_2$, resonance in $\mathbb R^3$, $A_2\times A_2$, generic $6$ in $\mathbb R^5$, non-generic $7$ in $\mathbb R^4$, $H_3$) & [AudCone]\\
$R_Pe_\zeta=0$ iff connected or $n\le2$; witness coefficient; $\psi_P$ an isomorphism; $\varphi_\Sigma$; $\Lie_ne_\zeta$; Prop.~4 maps; Theorem~$\Psi$ for random admissible $\lambda$ (\texttt{[wgt]}, fixed in \texttt{7c10288}, \S\ref{sec:natural}) &
\src{natural_isomorphism.py} & $n\le6$; $\psi_P$ and $R_Pe_\zeta$: also all $2045$ posets at $n=7$; $\varphi_\Sigma$ equivariance, Prop.~4: $n\le5$ & [Nat]\\
Independent audit of \S\ref{sec:natural}: Lemma~$\Psi$.0, contraction step, Theorem~$\Psi$ for every primitive $\zeta$ and random admissible $\lambda$, equivariance, Theorem~\ref{thm:ext}, $\Theta$ vs.\ extensions, $\varphi_\Sigma$; comb pairing $=\Sh_n$ without Ree &
\src{audit_natural_isomorphism.py} & all posets $n\le6$ (Lemma~$\Psi$.0, contraction, $\varphi_\Sigma$, equivariance, comb pairing: $n\le5$) & [AudNat]\\
\end{longtable}
}

\noindent Notes on the record:
\begin{itemize}[nosep]
\item No check in the table failed.
\item At $n=7$ the identity barely tests the refinement. For prime $n$,
  everything beyond Jen\v{c}a--Sarkoci and classical JKS is the vanishing
  $\mu(\OP(P)^g)=0$, and the sign factor is tested on a non-antichain poset
  only when $n$ is even and not a power of $2$ ([N7, \S2], \st{proven}). The
  $n=8,9,10$ families in [AudChar] were chosen to exercise the phases.
\item The $n=7$ run of \src{conjecture_s.py} is reported in [S] but was not
  archived, and the audit [AudS] did not re-run it (UNCHECKED).
\end{itemize}

\section{Open problems and unchecked items}\label{sec:open}

\begin{enumerate}
\item \textbf{Natural isomorphism: what remains.} The question is answered,
  in the precise sense of \S\ref{sec:natural}: an explicit, choice-free,
  relabelling-compatible isomorphism exists for every $P$
  (Theorem~\ref{thm:Theta}); naturality along extensions holds exactly on
  connected posets and is impossible at every disconnected $P$ with $n\ge3$
  (Theorem~\ref{thm:ext}); and compatibility with ordinal sums is either
  non-canonical or incompatible with the connected case. Still open:
  \begin{itemize}[nosep]
  \item Conjecture~\ref{conj:phisigma} ($\varphi_\Sigma$ invertible for
    every $P$);
  \item whether Theorem~\ref{thm:Psi} or Corollary~\ref{cor:klyachko}
    appears in the literature on Klyachko's theorem (one search only);
  \item a canonical $e_\zeta$-side coproduct of twisted form; every
    candidate twist found in [Nat, \S7] needs a base point;
  \item the \texttt{[wgt]} check of \src{natural_isomorphism.py} was
    defective; it is fixed in commit \texttt{7c10288}. The $n=7$ \texttt{[psi]} run, the
    $n=6$ and size-weighted $\varphi_\Sigma$ checks, the equivariance of
    $\varphi_\Sigma$ and the Dynkin formula were not re-run by [AudNat]
    (UNCHECKED); Props.~3--5 of [Nat] were audited by re-reading only.
  \end{itemize}
\item \textbf{Bracket, cooperad, Hopf-monoid structure.} The commutator of
  concatenations does \emph{not} descend to
  $\PLie(P)\otimes\PLie(Q)\to\PLie(P+Q)$ in the $e_\zeta$ model, even for
  antichains (\st{computed}, [Glob, \S3.4]). $\operatorname{Prim}(\mathbf Q)$
  carries a bracket, but it is not graded by posets. Open, and
  \st{conjectured} in [Glob, \S3.4, \S4.3]:
  \begin{itemize}[nosep]
  \item a poset cooperad on $\mathbb C[S_n]/\Sh_n$ via deconcatenation;
  \item a filtration of $\mathbf Q$ whose associated graded primitives are
    $\bigoplus_P\PLie(P)$, $P$-homogeneously;
  \item a ``local PBW'' $\mathbb C[\Lin P]\cong\operatorname{Sym}$ over
    components of the $\PLie(P|_T)$.
  \end{itemize}
\item \textbf{A $\Phi$-free proof of Theorem~\ref{thm:S}(c).} Run the
  induction of Theorem~\ref{thm:S}(a) directly on chains, by making precise
  the heuristic that $E_{a,x}$ is the $\Phi$-image of taking the link at the
  atom $\pi_{a,x}$ ([AudS, \S E.4]). This would remove the dependence on the
  cited Wachs statements for $n\ge7$.
\item \textbf{Cone/lune generalization.} [Cone, Conjecture~G]: the general
  form of Conjecture~S for every top-cone of every real central arrangement.
  The proven results of [Cone] have been independently audited in [AudCone]
  (the pointed-cone theorem holds for essential arrangements); Conjecture~G
  itself remains \st{conjectured}.
\item \textbf{Items the audits list as UNCHECKED.}
  \begin{itemize}[nosep]
  \item Checked in [Cit]: Hall (EC1 Prop.~3.8.5, in the form of
    Prop.~3.8.6); Weisner (EC1 Cor.~3.9.3, dual form); the Wachs 1998
    preprint numbers (Lemma~3.1, Thms~4.2--4.3, Prop.~2.3); the
    Jen\v{c}a--Sarkoci arXiv v2 numbering, including the identification of
    order congruences with acyclic-quotient partitions; and
    Cartier--Milnor--Moore/PBW for species (Aguiar--Mahajan, Ch.~11).
  \item NOT ACCESSIBLE ([Cit]): Crapo's theorem (Thm~3 according to
    Stanley, EC1 Ex.~3.92); the Hopf trace formula (Munkres Thm~22.1); Ree
    (Reutenauer Thm~3.1); and all journal or print numberings
    (Jen\v{c}a--Sarkoci in JCTA, Wachs in Discrete Math., EC1 in print,
    Aguiar--Mahajan \emph{Topics} in print; cf.\ [AudLit, \S5]).
  \item $\Sh_n\kappa_n\ne0$ for $n\ge6$, and hence the $P$-dependent
    $\kappa_n$-version of Prop.~C$'$ ([AudLit, rows 8, 30]).
  \item [AudLit, rows 24, 28, 29]: the conjectured restriction isomorphism
    of [Lit, \S2.3]; the labelling-dependent count of $W_Pe_\zeta=W_{P^{\rm op}}e_\zeta$;
    and $\dim W_P\kappa_n=e_C(P)$.
  \item The $n=7$ run of \src{conjecture_s.py} (not archived, not re-run).
  \item [Glob, \S\S3--4] (operations, $\mathbf Q$, primitives) has had no
    independent audit.
  \end{itemize}
\item \textbf{Novelty.} The literature search of [Lit, \S1] found no prior
  statement, proof or refutation of the conjecture, and no group action or
  Lie module attached to the homology of $\OP(P)$. Of the few papers citing
  Jen\v{c}a--Sarkoci, none studies a group action. [AudS, \S F] found no prior
  statement of Theorem~\ref{thm:S}(a)--(b) for general $P$. Neither search is
  exhaustive, and neither is a proof of novelty. Before publication, experts
  should be asked (e.g.\ on poset topology, free Lie algebras and
  arrangements) whether these results are folklore.
\item \textbf{The seed note's broader questions.} Its question about a
  ``primitive or Lie functor on posets'', and the moonshine-motivated search
  for an extra grading, are untouched.
\end{enumerate}

\begin{thebibliography}{99}

\bibitem{JencaSarkoci}
G.~Jen\v{c}a and P.~Sarkoci,
\emph{Linear extensions and order-preserving poset partitions},
J. Combin. Theory Ser.~A \textbf{122} (2014), 28--38,
doi:10.1016/j.jcta.2013.09.010; arXiv:1112.5782. Numbering used here is that
of arXiv v2.

\bibitem{Crapo}
H.~H.~Crapo, \emph{The M\"obius function of a lattice},
J. Combin. Theory \textbf{1} (1966), 126--131.

\bibitem{Wachs}
M.~L.~Wachs,
\emph{On the (co)homology of the partition lattice and the free Lie algebra},
Discrete Math. \textbf{193} (1998), 287--319. Theorem numbers cited here are
those of the preprint (Dec.~1994, revised Jan.~1996).

\bibitem{Ree}
R.~Ree, \emph{Lie elements and an algebra associated with shuffles},
Ann. of Math. \textbf{68} (1958), 210--220.

\bibitem{Reutenauer}
C.~Reutenauer, \emph{Free Lie Algebras}, LMS Monographs~7, Oxford, 1993.

\bibitem{Klyachko}
A.~A.~Klyachko, \emph{Lie elements in the tensor algebra},
Sibirsk. Mat. Zh. \textbf{15} (1974), no.~6, 1296--1304; English transl.
Siberian Math. J. \textbf{15} (1974), 914--920, doi:10.1007/BF00966559.

\bibitem{Stanley82}
R.~P.~Stanley, \emph{Some aspects of groups acting on finite posets},
J. Combin. Theory Ser.~A \textbf{32} (1982), 132--161.

\bibitem{Hanlon81}
P.~Hanlon, \emph{The fixed-point partition lattices},
Pacific J. Math. \textbf{96} (1981), 319--341.

\bibitem{Joyal86}
A.~Joyal, \emph{Foncteurs analytiques et esp\`eces de structures},
in: Combinatoire \'enum\'erative, LNM~1234, Springer, 1986, 126--159.

\bibitem{AguiarMahajan}
M.~Aguiar and S.~Mahajan, \emph{Topics in Hyperplane Arrangements},
Math. Surveys Monogr.~226, AMS, 2017; and \emph{Bimonoids for Hyperplane
Arrangements}, Encyclopedia Math. Appl.~173, CUP, 2020.

\end{thebibliography}

\noindent The bibliographic details of \cite{Stanley82,Hanlon81,Joyal86} are
as given in [Lit, \S5], marked there as ``from memory; verify''.

\end{document}
